%  LaTeX support: latex@mdpi.com 
%  For support, please attach all files needed for compiling as well as the log file, and specify your operating system, LaTeX version, and LaTeX editor.

%=================================================================
\documentclass[geomatics,article,accept,pdftex,oneauthor]{Definitions/mdpi} 

%=================================================================
%----------
% journal

% MDPI internal commands
\firstpage{1} 
\makeatletter 
\setcounter{page}{\@firstpage} 
\makeatother
\pubvolume{1}
\issuenum{1}
\articlenumber{0}
\pubyear{2025}
\copyrightyear{2025}
\externaleditor{Firstname Lastname}
\datereceived{9 December 2024} 
\daterevised{7 January 2025} 
\dateaccepted{15 January 2025} 
\datepublished{} 
\hreflink{\textls[-15]{https://doi.org/}} % If needed use \linebreak
%\doinum{}

%=================================================================
% Add packages and commands here. The following packages are loaded in our class file: fontenc, inputenc, calc, indentfirst, fancyhdr, graphicx, epstopdf, lastpage, ifthen, lineno, float, amsmath, setspace, enumitem, mathpazo, booktabs, titlesec, etoolbox, tabto, xcolor, soul, multirow, microtype, tikz, totcount, changepage, attrib, upgreek, cleveref, amsthm, hyphenat, natbib, hyperref, footmisc, url, geometry, newfloat, caption

\graphicspath{{figures/}} % path to folder with figures
%\usepackage{subcaption}
\usepackage{listings}
\usepackage{xcolor}
\usepackage{gensymb} % for \textdegree
\usepackage{tabularx}
\newcolumntype{Y}{>{\centering\arraybackslash}X}
\usepackage{amsmath,mathtools,amsthm}
	\setcounter{MaxMatrixCols}{15}
\usepackage{array}
\usepackage{amssymb}
\usepackage{amsfonts}
\usepackage{float}
\usepackage{chngcntr}
\counterwithin*{equation}{section}
\usepackage[super]{nth}
\usepackage{longtable}
\usepackage{tabularx}
\usepackage{multirow}
\usepackage{adjustbox}

\definecolor{deepblue}{rgb}{0,0,0.5}
\definecolor{deepred}{rgb}{0.6,0,0}
\definecolor{deepmagenta}{RGB}{195,12,214}
\definecolor{deepgreen}{RGB}{29,122,30}
\definecolor{mintgreen}{RGB}{192,249,192}
\definecolor{codepurple}{RGB}{204, 0, 153}
\definecolor{LightCyan}{rgb}{0.88,1,1}
\definecolor{GainsboroPurple}{RGB}{247,176,247}
\definecolor{LightSlateGrey}{RGB}{124,118,118}
%    
\lstdefinestyle{mystyle}{
    commentstyle=\color{LightSlateGrey},
    keywordstyle=\color{deepgreen}\bfseries,
    emphstyle=\color{deepred},
    numberstyle=\tiny\color{deepblue},
    stringstyle=\color{codepurple},
    basicstyle=\ttfamily\scriptsize,
    breakatwhitespace=false,         
    breaklines=true,                 
    captionpos=b,                    
    keepspaces=true,                 
    numbers=left,                    
    numbersep=5pt,                  
    showspaces=false,                
    showstringspaces=false,
    showtabs=false,                  
    tabsize=2
}
\lstset{style=mystyle}

%\usepackage{subcaption}
\usepackage{listings}
\usepackage{xcolor}
\usepackage{gensymb} % for \textdegree
\usepackage{amsmath,mathtools,amsthm}
	\setcounter{MaxMatrixCols}{15}
\usepackage{amssymb}
\usepackage{amsfonts}
\usepackage{float}
\usepackage{chngcntr}
\counterwithin*{equation}{section}
\usepackage[super]{nth}
\usepackage{longtable}
\usepackage{multirow}
\usepackage{adjustbox}
\usepackage{caption}
%\DeclareCaptionType{equ}[][]

%=================================================================
% Full title of the paper (Capitalized)
\Title{Improving Bimonthly Landscape Monitoring in Morocco, North Africa, by Integrating Machine Learning with GRASS GIS}

% MDPI internal command: Title for citation in the left column
\TitleCitation{Improving Bimonthly Landscape Monitoring in Morocco, North Africa, by Integrating Machine Learning with GRASS GIS}

% Author Orchid ID: enter ID or remove command
\newcommand{\orcidauthorA}{0000-0002-5759-1089} % Polina Add \orcidA{} behind the author's name

% Authors, for the paper (add full first names)
\Author{Polina Lemenkova $^{1,2}$\orcidA{}}
%MDPI: Please carefully check the accuracy of names and affiliations. Changes will not be possible after proofreading. %Author: yes, correct.


%\longauthorlist{yes}

% MDPI internal command: Authors, for metadata in PDF
\AuthorNames{Polina Lemenkova}

% MDPI internal command: Authors, for citation in the left column
\AuthorCitation{Lemenkova, P.}
%MDPI: Please check all author names carefully.

% If this is a Chicago style journal: Lastname, Firstname, Firstname Lastname, and Firstname Lastname.

% Affiliations / Addresses (Add [1] after \address if there is only one affiliation.)
\address{%
$^{1}$ \quad Department of Biological, Geological and Environmental Sciences, Alma Mater Studiorum---Università di Bologna, Via Irnerio 42, Emilia-Romagna, 40126 Bologna, Italy; polina.lemenkova2@unibo.it or polina.lemenkova@unibz.it; Tel.: +39-3446928732
%
$^{2}$ \quad Faculty of Agricultural, Environmental and Food Sciences, Libera Università di Bolzano, Piazza Università 5, Trentino-Alto Adige (S\"udtirol), 39100, Bolzano, Italy}
%MDPI:  affiliations: The address order should be from small to large. We have modified the address order. Please confirm. 
%Author: yes, confirm.
%MDPI: Alma Mater Studiorum---Università di Bologna/Trentino-Alto Adige|S\"udtirol: Please confirm that the address format is correct. %Author: yes, correct.
%MDPI: Please confirm that both email addresses are maintained. 
%Author: yes, confirm, I have two (2)  different emails from  two (2) universities.
% Contact information of the corresponding author
% Current address and/or shared authorship
%\firstnote{Current address: Affiliation 3.} 

% Abstract (Do not insert blank lines, i.e. \\) 
\abstract{This article presents the application of novel cartographic methods of vegetation mapping with a case study of the Rif Mountains, northern Morocco. The study area is notable for varied geomorphology and diverse landscapes. The methodology includes ML modules of GRASS GIS 'r.learn.train', 'r.learn.predict', and 'r.random' with algorithms of supervised classification implemented from the Scikit-Learn libraries of Python. This approach provides a platform for processing spatiotemporal data and satellite image analysis. The objective is to determine the robustness of the “DecisionTreeClassifier” and “ExtraTreesClassifier” classification algorithms. The time series of satellite images covering northern Morocco consists of six Landsat scenes for 2023 with a bimonthly time interval. Land cover maps are produced based on the processed, classified, and analyzed images. The results demonstrated seasonal changes in vegetation and land cover types. The validation was performed using a land cover dataset from the Food and Agriculture Organization (FAO). This study contributes to environmental monitoring in North Africa using ML algorithms of satellite image processing. Using RS data combined with the powerful functionality of the GRASS GIS and FAO-derived datasets, the topographic variability, moderate-scale habitat heterogeneity, and bimonthly distribution of land cover types of northern Morocco in 2023 have been assessed for the first time.}

% Keywords
\keyword{machine learning; remote sensing; image processing; satellite image; mapping; Africa; image analysis; automation; geography; earth sciences} 

% The fields PACS, MSC, and JEL may be left empty or commented out if not applicable
%\PACS{J0101}
%\MSC{}
%\JEL{}

\begin{document}

%----------- section ----------->
\section{Introduction}

The study of the relationship between geomorphology and vegetation has become a major issue for monitoring territories using Earth observation (EO) data \cite{Corenblit,Muir,Pu,Cary01091990}. The time series of such geospatial data are valuable for the evaluation of environmental changes based on their classification \cite{rs15225374,w17010068,app15010338}. Nevertheless, the issue of environmental monitoring using Remote Sensing (RS) data consists of receiving new data faster than we are able to describe, interpret, and process them. This is especially true in the satellite time series, where there is the need to classify, evaluate, and map habitats in operative regimes \cite{rs15225413,rs15225408,rs15225415}. In particular, geospatial data are used in the areas of agricultural land management and climate \mbox{change \cite{Ahmadian03052016,Lemenkova202412,IMTIAZ2024109172}}. The operative mapping of areas with diversified geomorphology is based on the time series of the satellite data that allow regular visualization for environmental monitoring \cite{JOHANSSON01011978,rs9121306}.  

In Morocco, the benefits from the use of satellite images for landscape analysis are made evident by its very diverse landscape structure \cite{Sahrane13122022,Marco04032021,Aouraghe01032010}. Since the geography of Morocco extends from the Atlantic Ocean to mountainous areas and the Sahara desert, this country is notable for its diversified geomorphological framework and varied topography. Moreover, the geological and tectonic development of the Atlas Mountains and their surroundings presents a complex pattern of geological units \cite{Benvenuti31122024,Ouallali03072020,Khalloufi01032010}. This also determines the distribution of major vegetation types, and controls land use patterns in the mountainous regions of the country \cite{ALIELOMAIRI2024101321}. 

In this context, land cover maps that visualize regional environmental settings play a vital role in mapping Morocco. Because RS data can be interpreted in terms of landscape structural units, comparing several images in a time series is useful for detecting changes. In turn, the change detection of land cover types enables major steps of landscape dynamics to be revealed \cite{CAI2023103226,Lemenkova202405,XIONG2022109132}. However, a question arises regarding the choice of classification algorithms and the efficiency of software for data processing. The integration of RS data with a Geographic Information System (GIS) for environmental mapping has been the subject of research for decades and continues to be a challenging topic. Many approaches to the problem of satellite image analysis have been proposed since the development of satellite missions \cite{NOY2023103377,Lemenkova202403,Su02062020,Potter18102014,Lemenkova202402}. A new line of research  arises at the intersection between programming methods of Machine Learning (ML) and spatial analysis using satellite images. The application of such tools and methods enables us to process images more effectively and accurately, as reported in previous studies \cite{rs16244812,w17010050,fire8010008}. 

Hence, the performance of image classification methods is linked to the quality of RS data. In the field of cartography, Landsat 8-9 Operational Land Imager and Thermal Infrared Sensor (OLI/TIRS) image data open new opportunities, since the spatial resolution of these data is an advantage for determining land use classes. In this context, the launch of new satellite constellations of Landsat 8-9 OLI/TIRS as EO systems allows the acquisition of time series. The acquisition of these data opens new perspectives in cartography \cite{Farah05012010,Lemenkova202322,Layati03042022,Layati28082024}.

%----------- section ----------->
\section{Objective and Motivation}

This paper focuses specifically on the application of ML methods to satellite image processing incorporated into the GRASS GIS software through the embedded Python Scikit-Learning libraries. Using ML in the traditional Geographic Information System (GIS) for mapping land cover types is a challenge that has been researched for many years and has many proposed solutions. Machine learning is an artificial approach to geospatial data processing that relies on mathematical and statistical algorithms to give computers the ability to “learn” from data. In this context, ML can be applied to satellite images such as Landsat 8-9 OLI/TIRS since pixels are discrete quantitative variables \cite{DAPAZGOMES,Lemenkova202026,LAROCHEPINEL2024109163}. 

Due to the specific advantages of ML, the functions of the programs improve their performance in solving geospatial analysis tasks \cite{More02072020,Lemenkova202403,GarciaPedrero03042017}. In the field of satellite image processing, this concerns the usual image processing tasks such as pixel recognition \cite{NASIRI2023103555,EBRAHIMY2023103390}, spectral reflectance analysis \cite{FU2024104421}, classification optimization \cite{PADHY2024100278,Lemenkova202401}, the development of more advanced methods of image analysis and interpretation, as well as the implementation of these methods to cartographic visualization, as discussed earlier \cite{HABIB20231,Lemenkova202406,KHALDI2024104191}. Compared to traditional approaches, the improvements brought by the ML algorithms concern the modeling of data with an average statistical error value % Attention AE: Check meaning retained.
 being as low as possible.

The Geographic Resources Analysis Support System (GRASS) GIS software was chosen as the satellite image processing tool for its combination of efficiency and cross-platform compatibility. This paper uses the GRASS GIS ML modules (`r.learn.train', `r.random', `r.learn.predict', and `r.category') with the `ExtraTreesClassifier' and `DecisionTreeClassifier' models for the supervised application, training, and classification of six satellite images. The results of the ML approach to image processing are compared with the results of unsupervised classification performed using the GRASS GIS modules `i.cluster' and `i.maxlik'. The scripts of this software are similar to libraries in programming languages.

In programming libraries, image processing algorithms are implemented as a set of functions used to process raster files, perform geospatial analysis, and create maps in the form of visualized images \cite{Lemenkova202227,PEREIRA2023106653,ESKANDARI2024}. This is the fundamental approach that differentiates the method of using script libraries from traditional GIS. The use of programming approaches to processing and analyzing satellite images allows the handling of complex representations of landscape features and detected objects recognized as land cover types that have various properties (extent and heterogeneity of landscapes, its texture, size of plots, fragmentation and topology, etc.).

As well as serving as the first descriptive study of habitat heterogeneity in the Rif Mountains, this study also provides guidance for future cartographic mapping using ML modules of GRASS GIS, taking into account the geographic isolation of diverse landscape patches and the heterogeneity and contrast of north African landscapes, ranging from arid areas of the Sahara to humid coastal zones. In particular, in this research, an exploration of GRASS GIS ML modules is presented and exploited, with the aim of demonstrating the applicability of spatiotemporal data processing to environmental studies in north Africa.

%----------- section ----------->
\section{Study Area}

The study focuses on northern Morocco, specifically the Rif Mountains; see Figure \ref{fig01}. The Rif Mountains region is located between latitudes 34\degree 30$^\prime$--35°15$^\prime$ N and longitudes 4°30$^\prime$--5°30$^\prime$ W, with an extension of about 30,000 m and an occupied area of 4000 km² \cite{ABOUTOFAIL2024112522}. The geographical position of the country and its vast extent result in a great variety of ecosystems. Landscapes vary from humid in the Rif and Atlas to the arid Saharan in the south, including sub-humid transition zones and a semi-arid climate in the plains and foothills areas \cite{MCGREGOR20091434}. 


The Rif Mountains were formed as a result of complex tectonic interactions between the African and Eurasian plates. The complex geological setting of northern Morocco in the Rif region has resulted in moderate tectonic activities \cite{CRINITI2024106861,PARISH199945}. 

The geomorphological structures of the region are composed of three main units from north to south: the internal zone, the flysch nappes, and the external zone \cite{ELTALIBI2014554,BONARDI2003149}. The crustal thickness beneath the Rif Mountains varies between 36 and 30 km towards the border with the western Mediterranean \cite{MAATE2000409}. The Rif presents a mountain range that extends across northern Morocco, skirting the Mediterranean Sea and in particular the Alboran Sea on its southern shore. The Rif constitutes a segment of the Alpine ranges. More specifically, it belongs to the Betico--Rifo--Tellian arc in the southwestern Mediterranean \cite{SLIMANI2019101785}. The Rif--Tell belt located in northern Morocco is a young Alpine orogenic system that presents metamorphic complexes formed by geodynamic and tectonic activity \cite{MARTINMARTIN2023106367,AITBRAHIM2002187,LEBLANC1984345}. This region is characterized by deep sediments related to Tethyan rifting \cite{ARANA1974197,PERRI2022105811,TALMAT2025105515}, which strongly affects soil properties and vegetation types \cite{CHAKKOUR2024,HAMOUCH2024}. As a result of such complex geologic and geomorphic settings in northern Morocco, there are very diverse ecosystem types that occupy the landscapes from the foothills, slopes, and valleys, from limestone mountains to the Saharan sebkhas. The main types of land cover are presented in Figure \ref{fig02}.

The geomorphological structure of the Rif can be divided into three segments. The Western Rif is covered with large mountains; its climate is mild, with heavy rains and showers in winter. The Central Rif has diverse reliefs and heterogeneous landscapes. The Eastern Rif is the flattest part of the range \cite{EDDAKIRI2024e02401}. The primary type is presented by the vegetation of the cultivated lands in mosaic, distributed along the basement of the Rif Mountains. The vegetation types in the Atlas Mountains region are presented by a closed to open evergreen or semi-deciduous broadleaf forest, with a forest cover of 100 to 40\% for closed to open trees or semi-deciduous trees \cite{BOUBEKRAOUI2023100388}. This geomorphological and climatic heterogeneity corresponds to bio-ecological and environmental diversity. Thus, the spatial richness of the conditions of climate and geology accentuates this biodiversity and the types of land use in Morocco.


\vspace{-9pt}
\begin{figure}[H]

	\hspace{-5pt}\includegraphics[width=13.5cm]{Fig_01}

    \caption{General topographic map of Morocco with outlined location of the study area. Mapping software: Generic Mapping Tools (GMT) version 6.4.0. Data source: General Bathymetric Chart of the Oceans (GEBCO). Map source: author.}
    \label{fig01}
\end{figure}

\vspace{-9pt}

\begin{figure}[H]

	\hspace{-5pt}\includegraphics[width=13.5cm]{Fig_02}

    \caption{Land cover types in Morocco. Data source: Food and Agriculture Organization (FAO). The map is created using QGIS software. Map source: author.}
    \label{fig02}
\end{figure}



The main structure of the Rif is represented by the limestone massifs. The underground water reserves of the Rif maintain the distribution of vegetation and the rivers that originate in the Rif \cite{Aqil03042015}. Due to these hydrogeological features, short coastal rivers and streams flow from the Rif to the Mediterranean, which are characterized by a torrential and erosional regime \cite{Zouhri01062004}. The sources of these rivers supply water to both diverse and very heterogeneous types of vegetation around the Rif. These processes briefly illustrate the close link between the geomorphological features, hydrogeology, and vegetation of the Rif region \cite{al2017milieux}. 

Seismic, magnetic, and thermal activities in the Rif caused by local crustal earthquakes are related to the geological instability of this region \cite{DUARTE2024104330}. Specifically, this region presents a southern branch of the Betic--Rif arc that is part of the Alpine ranges. This region has experienced deformations between northwest Africa and southern Spain (Iberia) since the Tertiary \cite{CHABLI2014123}, followed by the convergence between Africa and the Iberian \mbox{Peninsula \cite{LGHAMOUR2024108941,KAIROUANI2024106762}}. This tectonic block has moved towards the west-southwest and the region still experiences tectonic activity \cite{MALUSA2024230533}. At the same time, seismic-tectonic activities shape the regional geomorphology \cite{SADKAOUI2024105219}. In turn, this controls the distribution of land cover types and is expressed in the landscapes. Such links between geology, geomorphology, and the environmental context have been discussed in previous works \cite{COSTA2024109578,BOUTAHAR2023104828,KAIROUANI2024106762}.

The Mediterranean coasts of Morocco are a habitat for numerous wetland ecosystems, including peatlands with boreal floristic affinities and a high conservation value \cite{MULLER202235,AJBILOU2006104}. Occasional patches of rain-fed trees and scrub crops are distributed in the western regions of the country and depend on water availability \cite{MARANON1999147}. Rain-fed croplands, including olive trees and orchards in monoculture systems, interspersed with areas of sparse vegetation, are visible along the Atlantic Ocean coasts \cite{SALMON2015321,ZAYANI2023e22910}. Sparse herbaceous vegetation comprise complex patterns mixed with shrubs and thickets. Other land cover types are generally presented by urban areas around cities with artificial surfaces and associated areas and bare soil areas in desert and unoccupied lands; see Figure \ref{fig02}. 

%----------- section ----------->
\section{Materials and Methods}

%----------- subsection ----------->
\subsection{Data}
The Landsat 8-9 OLI/TIRS satellite image data covering the bimonthly period 2023: January, March, May, July, September, and November, Figure \ref{fig03}. 


{Although both Landsat 8 and Landsat 9 carry the Operational Land Imager (OLI) and the Thermal Infrared Sensor (TIRS), Landsat 9 is equipped with the second-generation versions of these instruments, known as the Operational Land Imager 2 (OLI-2) and the Thermal Infrared Sensor 2 (TIRS-2). In this study, images from Landsat 8 and 9 are employed. }The L2 data type is OLI\_TIRS\_L2SP with the ground control points of version 5. The processing software version is LPGS\_15.3.1c for Landsat-8 scenes and LPGS\_16.3.1 for Landsat-9 (2023) scene. The satellite roll angle is zero for Landsat 8 scenes and -0.001 for Landsat 9. The images were taken during the day at Nadir, with cloud cover less than 10\%. Here, only multispectral bands with an acquisition resolution of 30 m were used, except for panchromatic bands with a resolution of 15 m. The combination of multispectral bands is mainly used for the creation of false/true color composites for classification, especially for image visualization. The available data were listed using the search template in ``g.list rast''. 

\begin{figure}[H]
    \begin{center}
	\hspace{-15pt}\includegraphics[width=14.0cm]{Fig_03}
    \end{center}
    \caption{Landsat images on 2023: (\textbf{a})---January 26; (\textbf{b})---March 15; (\textbf{c})---May 10; (\textbf{d})---July 29; (\textbf{e})---September 23; (\textbf{f})---November 26.}
    \label{fig03}
\end{figure}



%----------- subsection ----------->
\subsection{Workflow}

We first plot the workflow obtained from each methodological step to provide a general picture of the approach used in this study; see Figure \ref{fig04}. 

Since various GISs include different steps of data processing, we next summarize the scheme into 12 well-recognized research steps, based on modules of GRASS GIS for geospatial data processing. Major research steps include the following points: (1) data collection; (2) software selection; (3) quality control; (4) data preprocessing; (5) image processing; (6) clustering (unsupervised classification); (7) machine learning; (8) training dataset; \mbox{(9) su}pervised classification; (10) numerical computation; (11) cartographic mapping; (12) validation. {The entire process of spatial RS data handling processing using GRASS GIS processes raw satellite images, extracts information from the original data, creates attributes obtained from the raster matrices, and enables use new information to digital maps. This links the original raw data to the final output of thematic environmental maps through the complex process of cartographic data management.}


\subsection{{Software}}
The distinct feature of GRASS GIS, the advanced software for RS data processing and cartographic mapping, is that {it} enables the performance of various tasks involved in cartographic tasks using a console-based menu. It also enables the integration of geospatial data, the extraction and interpretation of information from raw satellite images, the performance of mapping and visualization, the processing of numerical and statistical analysis and computation, and the evaluation of the results for quality control. Geospatial data that can be processed in GRASS GIS can originate from many different sources and it recognizes both raster and vector formats of diverse geospatial extensions. The software is freely and openly available through its developer portal 
(URL: \href{https://grass.osgeo.org/}{https://grass.osgeo.org/} (accessed on 8 December 2024)) which makes it a valuable tool for learning and research. The image processing process is composed of the following steps, implemented in the form of free and open source modules on GRASS GIS. The code for organizing the project and importing data is presented below. 
%MDPI: Please provide the access date of the URL in the following format: “URL (accessed on Day Month Year)”. 
%Author: ok, corrected.

\begin{figure}[H]
    \begin{center}
	\hspace{-35pt}\includegraphics[width=13.0cm]{Fig_04}
    \end{center}
    \caption{Workflow diagram illustrating major steps of research design used for the organization of data processing tasks. Software: GIMP, version 2.10. Diagram source: author.}
    \label{fig04}
\end{figure}
%MDPI: In principle, it is recommended that the figure should be placed after the first citation and not recommended to span the heading. However, we have spanned the third-level heading to reduce the space. Please confirm whether  keep. 
%Author: confirm in current location, agree with your corrections.

%----------- subsection ----------->


The GRASS GIS was created as the cartographic data program, and contains independent modules from diverse provider developers, for example, `r.import', `i.maxlik', etc. Moreover, GRASS GIS has a built-in framework and Python Application Programming Interface (API), which supports the advanced processing of the satellite time series analysis and spatial data. The mapping of topographic data was carried out with GMT using existing methods \cite{Lemenkova202221,BECKER20051075,Lemenkova202217,10537274,5369602}, while the mapping of land use data was carried out using QGIS \cite{FANG202031}. The mapping of land cover maps from RS data is based on satellite image classification methods.

%----------- subsection ----------->
\subsection{{Data Preprocessing}}
The data preprocessing procedures included several steps which include geometric and atmospheric correction for optical remote sensing images. First, the import and preprocessing of six Landsat images was implemented using the module `r.import' that reads the data. {Then, the} geometric correction was performed using the 'i.rectify' module of GRASS GIS which corrected each image by performing a coordinate transformation for the pixels in the images using control points. Hence, the x,y cell coordinates on the images were re-calculated to a transformation matrix and then converted to standard map coordinates for each pixel in the image. Following that, the atmospheric correction was performed using the 6S algorithm embedded in the 'i.atcorr' module of GRASS GIS. 

The atmospheric correction is essential to numerically evaluating the amount of wavelengths reflected in the diverse land cover types of the surface in various seasonal \mbox{periods \cite{f14071414,ABDELRAHMANHEGAB2024716}}. {The} `i.landsat.toar' was used to convert Landsat 8--9 OLI/TIRS images and to transform the calibrated digital numbers for the reflectance at the top of the atmosphere. In this way, the quality of the images has been improved. The pansharpening has been performed using the `i.pansharpen' module of GRSAS GIS which uses image fusion algorithms to sharpen multispectral channels with high-resolution panchromatic channels. In this way, the resolution of the multispectral image is improved. The effectiveness of the pansharpening is described in earlier studies on RS data processing \cite{rs15082017,rs13173433}.

%----------- subsection ----------->
\subsection{{Creating Color Composites}}

The creation of color composites (Figure \ref{fig05}) was performed based on the Landsat multispectral bands using the `r.composite' module and the visualization of the images using the `d.rast' module. The spectral combination calculation integrates different spectral bands of Landsat based on the physical properties of the land covers reflected in the images. Therefore, the spectral bands are useful to improve the discrimination between landscape classes. The estimation of the color composition model aims to better recognize the ten land cover classes, representing landscape patches in the Rif Mountains and its surroundings in the images. The vegetation observed on the Atlantic Ocean side appears in darker reds than the surrounding vegetation around the Rif Mountains. The water that absorbs almost all wavelengths is very dark, up to black, while the empty ground surface appears very light, in tones ranging from beige to ochre; see Figure \ref{fig05}. On the colored composition in ``false colors'', the vegetation appears in different shades depending on the species, but also on the environmental conditions: the geomorphology of the mountain slope, the curvature, and the inclination of the slope affect the shadows on the images; Figure \ref{fig05}.

\begin{figure}[H]

\includegraphics[width=13.5cm]{Fig_05}

    \caption{Landsat 8-9 multispectral images: false and true color composites.}
    \label{fig05}
\end{figure}

\begin{footnotesize}
\begin{verbatim}
# Create the project with new location from raster map 
(file must contain projection metadata):
# grass -c myraster.tif /home/user/grassdata/mynewlocation
grass
#cd /Users/polinalemenkova/grassdata
#grass -c LC09_L2SP_179073_20220419_20230421_02_T1_SR_B1.tif 
/Users/polinalemenkova/grassdata/Morocco
# ----IMPORT AND PREPROCESSING-------------------------->
# g.mapset location=Morocco mapset=PERMANENT
g.list rast
# importing the image subset with 7 Landsat bands and display the raster map
r.import input=/Users/polinalemenkova/grassdata/Morocco/
LC09_L2SP_201036_20230126_20230313_02_T1_SR_B1.TIF 
output=L9_2023_J_01 extent=region resolution=region --overwrite
r.import input=/Users/polinalemenkova/grassdata/Morocco/
LC09_L2SP_201036_20230126_20230313_02_T1_SR_B2.TIF 
output=L9_2023_J_02 extent=region resolution=region
r.import input=/Users/polinalemenkova/grassdata/Morocco/
LC09_L2SP_201036_20230126_20230313_02_T1_SR_B3.TIF 
output=L9_2023_J_03 extent=region resolution=region
r.import input=/Users/polinalemenkova/grassdata/Morocco/
LC09_L2SP_201036_20230126_20230313_02_T1_SR_B4.TIF
 output=L9_2023_J_04 extent=region resolution=region
r.import input=/Users/polinalemenkova/grassdata/Morocco/
LC09_L2SP_201036_20230126_20230313_02_T1_SR_B5.TIF
 output=L9_2023_J_05 extent=region resolution=region
r.import input=/Users/polinalemenkova/grassdata/Morocco/
LC09_L2SP_201036_20230126_20230313_02_T1_SR_B6.TIF 
output=L9_2023_J_06 extent=region resolution=region
r.import input=/Users/polinalemenkova/grassdata/Morocco/
LC09_L2SP_201036_20230126_20230313_02_T1_SR_B7.TIF
 output=L9_2023_J_07 extent=region resolution=region
g.list rast
\end{verbatim}
\end{footnotesize}

The modeling and simulation of the images using the principle of color compositions includes `true color' and `false color' images; see Figure \ref{fig05}. Here, the spectral bands of the images are visualized in several color combinations that are widely used in RS because they are very suitable for studying vegetation and landscape. For example, the combination with band number 5 (infrared, IR) is based on the properties of vegetation that reflect very strongly near-IR radiation.

\begin{footnotesize}
\begin{verbatim}
# ----CREATING COLOR COMPOSITES----->
# false color
r.composite blue=L9_2023_J_07 green=L9_2023_J_05 red=L9_2023_J_03
 output=L9_2023_J_753 --overwrite
d.mon wx0
d.rast L9_2023_J_753
d.out.file output=Morocco_753 format=jpg --overwrite
# false color: NIR band B05 in the red channel, red band B04 in the green 
channel and green band B03 in the blue channel
r.composite blue=L9_2023_J_03 green=L9_2023_J_04 red=L9_2023_J_05
 output=L9_2023_J_345 --overwrite
d.mon wx0
d.rast L9_2023_J_345
d.out.file output=Morocco_345 format=jpg --overwrite
# true color
r.composite blue=L9_2023_J_02 green=L9_2023_J_03 red=L9_2023_J_04
output=L9_2023_J_234 --overwrite
d.mon wx0
d.rast L9_2023_J_234
d.out.file output=Morocco_J_234 format=jpg --overwrite
\end{verbatim}
\end{footnotesize}

%----------- subsection ----------->
\subsection{{Image Classification}}
The aim of {image classification} is to model the {raster }data and predict for each pixel a label, i.e. a Land Cover Class. {In this work, image classification was performed using} two {principally distinct approaches}: supervised classification and unsupervised classification. The use of an ML approach for satellite image processing by GRASS GIS consists of two phases. The first {includes} estimating a model from data by automatic clustering. The unsupervised clustering approach groups comparable samples within the same class using the k-means algorithm. The image clustering and classification groups samples according to the maximum likelihood algorithm (unsupervised learning) using the `i.cluster' and `i.maxlik' modules. The groups, i.e., the clusters, are made up of pixels with similar spectral values. 

The unsupervised classification employs a clustering approach performed using the ``i.cluster'' module. The ``i.cluster'' generates a signature file and reports the cluster maps, or regions, that were automatically generated using the ``k-means'' data partitioning algorithm. The GRASS SIG code used for this is the following (example for the 2013 image, repeated for images from all other years using the same technique): ``i.cluster group=L8\_2013 subgroup=res\_30m signaturefile=cluster\_L8\_2013 classes=8 reportfile=rep\_clust\_L8\_2013.txt''. Therefore, samples within a class are dissimilar from the samples in others.

\begin{footnotesize}
\begin{verbatim}
# ---Clustering and Classification ---->
# grouping data by i.group
# Set computational region to match the scene
g.region raster=L9_2023_J_01 -p
i.group group=L9_2023_J subgroup=res_30m \
input=L9_2023_J_01,L9_2023_J_02,L9_2023_J_03,L9_2023_J_04,
L9_2023_J_05,L9_2023_J_06,L9_2023_J_07 --overwrite
# Clustering: generating signature file and report using 
k-means clustering algorithm
i.cluster group=L9_2023_J subgroup=res_30m \
  signaturefile=cluster_L9_2023_J \
  classes=10 reportfile=rep_clust_L9_2023_J.txt --overwrite
# Classification by i.maxlik module
i.maxlik group=L9_2023_J subgroup=res_30m \
  signaturefile=cluster_L9_2023_J \
  output=L9_2023_J_cluster_classes reject=L9_2023_J_cluster_reject 
# Mapping
d.mon wx0
g.region raster=L9_2023_J_cluster_classes -p
r.colors L9_2023_J_cluster_classes color=bcyr
d.rast L9_2023_J_cluster_classes
d.legend raster=L9_2023_J_cluster_classes title="26 January 2023" 
title_fontsize=14 font="Helvetica" fontsize=12 bgcolor=white
 border_color=white
d.out.file output=Morocco_2023_Jan format=jpg --overwrite
# Mapping rejection probability
d.mon wx1
g.region raster=L9_2023_J_cluster_classes -p
r.colors L9_2023_J_cluster_reject color=wave -e
d.rast L9_2023_J_cluster_reject
d.legend raster=L9_2023_J_cluster_reject title="26 January 2023" 
title_fontsize=14 font="Helvetica" fontsize=12 bgcolor=white 
border_color=white
d.out.file output=Morocco_2023_reject format=jpg --overwrite
\end{verbatim}
\end{footnotesize}

In this way, each land cover class is associated with each cluster, correspondently. Calculating the pixel rejection probability was performed using accuracy assessment for quality control. After the unsupervised modeling is finished, the new data obtained are submitted as the training dataset for ML to obtain better image classification results. The clustering results are used as observation data available for supervised learning. The results of this step are used for ML as the training dataset. Subsequently, the ML processing of images was performed using supervised classification that uses a sequence of modules. The ability of ML to model geospatial data relies heavily on the acquisition of training data as the algorithm feeds on the input dataset, so this is an important step. ML classifies the images % Attention AE: Check meaning retained.
using the “r.learn.train” and “r.learn.predict” models.

This algorithm is based on the use of the ‘decision tree’ and ‘extra trees’ classifiers that are integrated into the GRASS GIS module ‘r.learn.train’. Here, the data are labeled, i.e., this is supervised learning with the technical specificities mentioned in the code below. First, the ‘r.random’ module was used to generate training pixels from a previous classification of the land cover of Morocco. All these algorithms were implemented using the GRASS GIS module ‘r.learn.train’ using the programming code presented below.

\begin{footnotesize}
\begin{verbatim}
# MACHINE LEARNING ------>
# g.list rast
g.region raster=L9_2023_J_01 -p
# First, we are going to generate some training pixels from an older (1996) 
land cover classification:
r.random input=L9_2023_cluster_classes seed=100 npoints=1000 
raster=training_pixels --overwrite
# Next, we create the imagery group with all 
Landsat-8 OLI/TIRS 7 (2000) bands:
i.group group=L9_2023_J 
input=L9_2023_J_01,L9_2023_J_02,L9_2023_J_03,L9_2023_J_04,
L9_2023_J_05,L9_2023_J_06,L9_2023_J_07 --overwrite
# Then use these training pixels to perform a classification on 
recent Landsat - 2022 image:
#
# train a decision tree classification model using r.learn.train
r.learn.train group=L9_2023_J training_map=training_pixels \
    model_name=DecisionTreeClassifier n_estimators=500 
    save_model=rf_model.gz --overwrite
# perform prediction using r.learn.predict
r.learn.predict group=L9_2023_J load_model=rf_model.gz
 output=rf_classification --overwrite
# check raster categories - they are automatically applied 
to the classification output
r.category rf_classification
# copy color scheme from landclass training map to result
r.colors rf_classification raster=training_pixels
# display
d.mon wx0
d.rast rf_classification
r.colors rf_classification color=roygbiv -e
d.legend raster=rf_classification title="Decision Tree: 01/2023" 
title_fontsize=14 font="Helvetica" fontsize=12 bgcolor=white 
border_color=white
d.out.file output=DT_2023_01 format=jpg --overwrite
\end{verbatim}
\end{footnotesize}

In the next step, the prediction was performed using the module ``r.learn.predict'' which was performed in the code: ``r.learn.predict group=L8\_2013 load\_model=rf\_model.gz output=rf\_classification''. The raster categories were examined using the module ``r.category'' of GRASS SIG, based on the data automatically applied to the classification output. For this, the following code was used: ``r.category rf\_classification''. The maps were visualized using the following combination of two main modules: ``d.rast'' which visualizes the map itself and ``d.legend" which adds the legend to this map. The images were converted to bitmap format using the module ``d.out.file", e.g., ``d.out.file output=Maroc\_2013 format=jpg''.

\begin{footnotesize}
\begin{verbatim}
# train a extra trees classification model using r.learn.train
r.learn.train group=L9_2023_J training_map=training_pixels \
   model_name=ExtraTreesClassifier n_estimators=500 
   save_model=rf_model.gz --overwrite
# perform prediction using r.learn.predict
r.learn.predict group=L9_2023_J load_model=rf_model.gz  output=rf_classification --overwrite
# check raster categories applied to the classification output
r.category rf_classification
# copy color scheme from landclass training map to result
r.colors rf_classification raster=training_pixels
# display
d.mon wx0
d.rast rf_classification
r.colors rf_classification color=bgyr -e
d.legend raster=rf_classification title="Extra Tree: 01/2023" 
title_fontsize=14 font="Helvetica" fontsize=12 bgcolor=white 
border_color=white
d.out.file output=ET_2023_01 format=jpg --overwrite
\end{verbatim}
\end{footnotesize}

This was used as a basis for ML using the following code: “r.random input=L8\_2013\_cl\linebreak  uster\_classes seed=100 npoints=1000 raster=training\_pixels”. After that, the “r.learn.train” module was used to train a module using the training map that was created in the previous step by applying the built-in algorithm. The code used is as follows: “r.learn.train group=L8\_2013 training\_map=t\_pixels model\_name=DecisionTreeClassifier n\_estimators=500 save\_model=rf\_model.gz”. For all approaches, the model name was changed using the “model\_name” function and the name of the model in question, e.g., “modelage=RandomForestRegressor”.

Since the pixel labels are discrete (and not continuous which requires the use of regression), the classification methods `decision trees' and `supplementary trees' were chosen. These ML algorithms predict the probability of pixel values suitable for the target class on the image. Therefore, the machine performs a classification of the image pixels into 10 land cover classes on the Rif Mountains. The categorization of the classes was carried out using the `r.category' module of GRASS GIS. 

Finally, the visualization of images was achieved using the auxiliary graphic modules of GRASS GIS for cartographic processing: `d.mon', 'd.rast', `r.colors' and `d.legend'. The images were saved as bitmap files using the module `d.out.file'. The calculation and estimation of classification accuracy are carried out by the function `reject' for the evaluation of the probability of the percentage of correctly classified pixels on the image. The evaluation of image processing quality and algorithm performance is conducted by GRASS GIS.

%----------- section ----------->
\section{Results and Discussion}

%----------- subsection ----------->
\subsection{{Analysis of Findings}}

Mountain areas in northern Morocco are particularly sensitive to climate change and pastoral activities. Seasonal variations in climate and effects from human activities strongly influence agro-ecosystems and cause changes in land cover types. Although there are currently some efforts to map the landscape dynamics in such areas, any seasonal response of the land cover types will be better evaluated using a quantitative assessment of the RS data and ML methods. Using examples from the Rif Atlas in Morocco, this paper explored the seasonal variations in land cover types in northern Morocco in 2023 caused by the effects from livestock management, arboriculture, and climate change. The selected period was assessed for the year 2023 using the available time series from Landsat dataset collection. The results of satellite image classification were obtained using three % Attention AE: Confirm whether should be two.
approaches: unsupervised classification and supervised classification by ML methods which employs two distinct algorithms: ``Decision tree classifier'' and ``Extra tree classifier''. Figure \ref{fig06} shows the images generated using the clustering classification method, which illustrate the land cover changes with a bi-monthly time interval; see Table \ref{tab01}.

\begin{figure}[H]

	\hspace{-0pt}\includegraphics[width=13.5cm]{Fig_06}

    \caption{Satellite image classification using unsupervised classification clustering methods and the GRASS GIS `MaxLike' approach.}
    \label{fig06}
\end{figure}


\vspace{-5pt}
\begin{table}[H]
\caption{Bimonthly variations in Rif Mountains by pixels over 10 land cover classes in 2023 in two (cold/hot) different seasons: (1) January, March, and May; (2) July, September, and November.}
\begin{tabularx}{\textwidth}{p{1.5cm}XXX X X  X  XXX X  }
  \toprule
   \textbf{Class} & \textbf{1} & \textbf{2} & \textbf{3} & \textbf{4} & \textbf{5} & \textbf{6} & \textbf{7} & \textbf{8} & \textbf{9} & \textbf{10} \\
   \midrule
    January & 766 & 465 & 554 & 736 & 1031 & 785 & 462 & 1016 & 646 & 149 \\
    March & 720 & 341 & 696 & 673 & 1060 & 791 & 587 & 821 & 720 & 203 \\
    May & 742 & 348 & 545 & 449 & 904 & 912 & 905 & 722 & 645 & 501 \\
    July & 766 & 465 & 554 & 736 & 1031 & 785 & 462 & 1016 & 646 & 149 \\
    September & 668 & 424 & 674 & 318 & 793 & 791 & 1099 & 1023 & 683 & 222 \\
    November & 781 & 483 & 701 & 675 & 1097 & 429 & 744 & 833 & 734 & 225 \\
    \bottomrule
\end{tabularx}
\label{tab01}
\end{table}

%----------- subsection ----------->
\subsection{{Interpretation of Land Cover Types}}

Figure \ref{fig07} shows that the classification using the “decision tree classifier” method improved the results since the ML results are more robust and stable compared to the unsupervised classification. The January scene shows the least plant activity which gradually increases from March to May and decreases, respectively, after the hot summer period until November. {Vegetation} varied seasonally in the Rif region according to the natural dynamics of the phenological cycle during different calendar months. {The maps show seasonal changes in 10 categories of land cover classes that reflect landscape variations: (1) water; (2) broadleaved forest; (3) grassland; (4) wetlands; (5) urban areas; (6) coniferous forest; (7) shrubland (thicket); (8) sparse vegetation in semi-deserts; (9) bare areas; (10) cropland. }The comparison of the obtained results was performed using the existing classification schemes in studies on land cover types in Morocco \cite{Kusi,Hamdani,Kusi2023,Elfallah,Idoumskine}.

The identified 10 land cover classes were generalized from the FAO-based land cover classes in Morocco; see Figure \ref{fig02}. {The most important type of land cover categories in northern Morocco is pastures, which include steppes, high meadows, and semi-desert lands with occasional vegetation inclusions. Pastures provide nourishment for livestock and resources for population. Nevertheless, as a result of overgrazing without land management, millions of hectares of pastures are being degraded in northern Morocco which causes environmental degradation. On the satellite images, the decrease in land cover types `Grassland' and 'sparse vegetation in semi-deserts' is notable through the comparison of the classified images.} 

\begin{figure}[H]
    \begin{center}
	\hspace{-0pt}\includegraphics[width=13.5cm]{Fig_07}
    \end{center}
    \caption{Results of satellite image classification using supervised ML method with GRASS GIS `Decision Tree' algorithm.}
    \label{fig07}
\end{figure}

{Another environmental issue in Morocco is presented by the degradation of forests, which is mostly caused by human activities such as resource exploration (quarries and other geologic activities), urbanization with the associated extension of built-up areas and constructed roads, the effects from agricultural activities which include the replacement of areas formerly occupied by forest to croplands and agricultural lands, and finally, the overgrazing and conversion of natural lands to rangeland. Deforestation is visible by comparing classified satellite images (class ``Forests'').}

Compared to other land cover types, the mosaic vegetation types of deciduous forests developed partially and differently for both algorithms, with more density and a homogeneous pattern across the entire study area. Figure \ref{fig08} shows the land cover type patterns over the Rif Mountains and surrounding areas, defined using an improved ML algorithm. The analysis of the presented maps approved the previously reported current environmental problems of Morocco which include deforestation, overgrazing \cite{ZAHER2020109544}, soil erosion, landscape deterioration and desertification \cite{MOUMANE2022100745,MULAS2012e46,KOUBA20181}, water scarcity, and climate change \cite{MAYAUX2024104093}. Besides natural variability of the Rif region, forest areas are affected by human factors (socio-demographic, economic, and urban dynamics) and hazards caused by erosion and forest fires \cite{chebli2018forest}. The expansion of agricultural lands negatively affects the mountainous forests of the Rif \cite{Adghir02112018,Boutahar01102024,Jimenez01092002}. In addition, the contrasts from hot and humid climates and the varied topography of the Rif facilitates landslides that present geomorphological risks in Northern Morocco \cite{Dahmani10102024,EsSmairi13122022,Manaouch03032024}.

\begin{figure}[H]

	\hspace{-0pt}\includegraphics[width=13.7cm]{Fig_08}

    \caption{Results of satellite image classification using supervised ML method with GRASS GIS `Extra Tree' algorithm.}
    \label{fig08}
\end{figure}

%----------- subsection ----------->
\subsection{{Comparison Between Performance of Algorithms}}

{Selecting the optimal method of land cover modeling ensures the formalization of our understanding of Earth's environmental system and dynamics through the quantification of various landscape elements. By quantifying the distribution of land cover types, cartographic models enable the interpretation of relations between patches that constitute a complex mosaic of Earth's landscape which integrates different components into a coherent framework. Therefore, technical} improvements of the parameters that differ in various algorithms ‘MaxLike’, ‘DecisionTreeClassifier’, and ‘ExtraTreesClassifier’ {were assessed to analyse the best performance of the algorithms for} land cover mapping. 

{The} most complex landscape patches with high heterogeneity can be better distinguished, including sparsely distributed and densely distributed vegetation segments. Compared to the previous results in Figures \ref{fig06} and \ref{fig07}, the results of the Extra tree are not quite similar for the classification maps based on the “Decision tree” (Figure \ref{fig07}) and “Extra tree” algorithms (Figure \ref{fig08}), because they rely on the extremely random trees method. The “Extra tree” algorithm is similar to the random forest algorithm in its nature, but it processes images much faster. Therefore, technically, this algorithm is an improved ensemble supervised ML method.

However, this algorithm creates many decision trees using random sampling of pixels for each tree. The most important and unique feature of the “Extra trees” algorithm is that it randomly selects a pixel division value. Therefore, compared with the “decision tree” algorithms, the land cover types are diverse and uncorrelated. The best performance among the three image processing methods is demonstrated by the “decision tree” algorithms, which divide the image into land cover classes with precision and detail. The accuracy assessment is presented in Figure \ref{fig09} and the computational results are available in Tables \ref{tab03} and \ref{tab04}.




\begin{table}[H]


\caption{Bimonthly evaluation of the class means of Digital Numbers (DN) computed for pixels assigned to 10 land cover classes in landscapes of Rif Mountains, north Morocco. Each multispectral band (1 to 7) of Landsat 8-9 OLI/TIRS is evaluated. Period: 2023, January to May.} \label{tab03}

\begin{tabularx}{\textwidth}{@{\extracolsep{\fill}}p{1.5cm}cccccccc@{}}

\toprule  \textbf{Class} & \textbf{Band 1} & \textbf{Band 2} & \textbf{Band 3} & \textbf{Band 4} & \textbf{Band 5} & \textbf{Band 6} & \textbf{Band 7} \\  
\midrule

%\bf{2019} &  &  &  &  &  &  & \\
\multicolumn{8}{c}{\textbf{2023: January}}\\
\midrule
1 & 8053.66 & 8105.74 & 7961.36 & 7451.51 & 7386.15 & 7411.52 & 7385.87 \\
2 & 7529.44 & 7754.46 & 8346.18 & 8318.72 & 11,431.1 & 10,065.3 & 8996.52 \\
3 & 8054.51 & 8324.72 & 9136.1 & 9317.05 & 12,308.5 & 12,725.3 & 11,052.1 \\
4 & 7889.8 & 8102.02 & 8942.29 & 8745.57 & 14,901.1 & 12,282.7 & 10,135 \\
5 & 8332 & 8667.33 & 9728.08 & 10061.3 & 14,289.8 & 14,818.8 & 12,637.9 \\
6 & 8110.13 & 8324.27 & 9467.21 & 9012.78 & 17,982.6& 13,440.9 & 10,707.4 \\
7 & 8251.2 & 8447.13 & 9831.94 & 9044.77 & 21,802.8 & 13,717.8 & 10,608.9 \\
8 & 8640.53 & 9039.94 & 10,407.5 & 10,794.4 & 16,840.5 & 16,286.3 & 13,653.7 \\
9 & 9111.71 & 9625.79 & 11,316.9 & 12,294.1 & 17,724.7 & 19,083.2 & 16,430 \\
10 & 11,270 & 12,124 & 14,516 & 16,218.7 & 21,505.9 & 23,336.4 & 20,483 \\
\midrule
%\bf{2022} &  &  &  & & & & \\
\multicolumn{8}{c }{\textbf{2023: March}}\\
\midrule
1 & 7732.78 & 7928.37 & 7997.69 & 7540.35 & 7530.62 & 7850.57 & 7839.92 \\
2 & 7850.68 & 8119.44 & 9096.92 & 8986.22 & 16,313.3 & 12,929.6 & 10,370.1 \\
3 & 8406.34 & 8821.78 & 10,011.9 & 10,547.6 & 15,377.5 & 14,966.2 & 12,509.7 \\
4 & 8377.33 & 8825.2 & 10,246.8 & 10,459 & 19,450.9 & 15,628.6 & 12,195.9 \\
5 & 8777.26 & 9318.35 & 10,778.5 & 11,790.6 & 16,488.5 & 17,202.7 & 14,298.8 \\
6 & 9050.82 & 9764.11 & 11,495.9 & 12,894.6 & 18,940.7 & 17,954 & 14,120.4 \\
7 & 9129.74 & 9761.09 & 11,407.2 & 12,937.6 & 17,616.1 & 19,787.9 & 16,566.3 \\
8 & 9848.06 & 10,823.1 & 12,951.4 & 15,075.9 & 19,854.8 & 20,511.1 & 16,312.7 \\
9 & 9570.95 & 10,302.4 & 12,240.4 & 14,420.2 & 19,123.6 & 22,705 & 19,360.4 \\
10 & 10,759 & 11,957.4 & 14,571.2 & 17,344.9 & 21,903.1 & 24,000.4 & 19,583.5 \\
\midrule

%\bf{2023} &  &  &  & & & & \\
\multicolumn{8}{ c }{\textbf{2023: May}}\\
\midrule
1 & 8111.51 & 8199.08 & 8100.88 & 7465.42 & 7287.35 & 7335.55 & 7330.84 \\
2 & 7565.9 & 7794.02 & 8480.23 & 8399.37 & 13,105.3 & 11,135.9 & 9540.2 \\
3 & 7941.53 & 8209.27 & 9135.12 & 9131.33 & 14,663.5 & 13,138.1 & 10,966.2 \\
4 & 8278.6 & 8635.24 & 9757.47 & 10183.6 & 14,710.6 & 15,182 & 12,997.8 \\
5 & 8007.61 & 8238.4 & 9379.12 & 9007.12 & 18,137.6 & 13,349.7 & 10,649.6 \\
6 & 7959.3 & 8113.15 & 9263.84 & 8488.51 & 22,413.9 & 12,764.6 & 9831.97 \\
7 & 8460.35 & 8799.31 & 10,364.7 & 10,092.9 & 20,827.9 & 15,638.7 & 12,351.6 \\
8 & 8552.87 & 8939.44 & 10,306.6 & 10,660.9 & 17,349.2 & 16,382.3 & 13,624.4 \\
9 & 8999.75 & 9517.87 & 11,173.8 & 12,159.9 & 17,889.8 & 18,838.8 & 16,259.8 \\
10 & 10,309.3 & 11,076.4 & 13,413.5 & 15,174.4 & 20,565.5 & 22,691.1 & 20,127.9 \\
\bottomrule
\end{tabularx}
\end{table}
%MDPI: Please confirm if the bold is unnecessary and can be removed. The following highlights are the same.
%Author: bold is not necessary, I removed it.


\begin{adjustwidth}{-\extralength}{0cm}
\begin{equation}
\scriptsize
\begin{vmatrix}
   & 1 & 2 & 3 & 4 & 5 & 6 & 7 & 8 & 9 & 10 \\
1 & 0 &  &  &  &  &  &  &  &  &  \\
2 & 1.7 & 0 &  &  &  &  &  &  &  &  \\
3 & 3.3 & 1.1 & 0 &  &  &  &  &  &  &  \\
4 & 3.3 & 1.0 & 0.7 & 0 &  &  &  &  &  &  \\
5 & 4.2 & 1.9 & 0.9 & 1.2 & 0 &  &  &  &  &  \\
6 & 3.9 & 1.6 & 1.2 & 0.8 & 1.0 & 0 &  &  &  &  \\  
7 & 3.5 & 1.8 & 1.6 & 1.2 & 1.2 & 0.6 & 0 &  &  &  \\  
8 & 4.6 & 2.3 & 1.5 & 1.7 & 0.7 & 1.2 & 1.1 & 0 &  &  \\
9 & 5.0 & 3.0 & 2.2 & 2.5 & 1.5 & 1.9 & 1.6 & 0.9 & 0 &  \\   
10 & 3.5 & 2.6 & 2.2 & 2.3 & 1.8 & 2.0 & 1.7 & 1.4 & 0.9 & 0 \\          
\end{vmatrix}
 %
\hspace{3pt}
\begin{vmatrix}
    & 1 & 2 & 3 & 4 & 5 & 6 & 7 & 8 & 9 & 10 \\
1 & 0 &  &  &  &  &  &  &  &  &  \\
2 & 2.5 & 0 &  &  &  &  &  &  &  &  \\
3 & 3.0 & 1.1 & 0 &  &  &  &  &  &  &  \\
4 & 3.5 & 1.3 & 0.6 & 0 &  &  &  &  &  &  \\
5 & 4.3 & 2.2 & 0.9 & 0.9 & 0 &  &  &  &  &  \\
6 & 4.7 & 2.6 & 1.3 & 1.1 & 0.7 & 0 &  &  &  &  \\  
7 & 5.1 & 3.1 & 1.7 & 1.7 & 1.0 & 0.8 & 0 &  &  &  \\  
8 & 5.4 & 3.8 & 2.3 & 2.3 & 1.7 & 1.1 & 1.0 & 0 &  &  \\
9 & 5.7 & 3.8 & 2.5 & 2.4 & 1.9 & 1.5 & 1.0 & 0.8 & 0 &  \\   
10 & 5.5 & 4.2 & 3.0 & 2.9 & 2.6 & 2.0 & 1.9 & 1.1 & 1.1 & 0 \\              
\end{vmatrix} 
\end{equation}
\end{adjustwidth}
1: Bimonthly class separability matrices between 10 classes of land cover types in Rif Mountains, Morocco: January and March 2023 (1).


\begin{adjustwidth}{-\extralength}{0cm}
\begin{equation}
\scriptsize
\begin{vmatrix}
   & 1 & 2 & 3 & 4 & 5 & 6 & 7 & 8 & 9 & 10 \\
1 & 0 &  &  &  &  &  &  &  &  &  \\
2 & 2.6 & 0 &  &  &  &  &  &  &  &  \\
3 & 4.0 & 0.9 & 0 &  &  &  &  &  &  &  \\
4 & 4.5 & 1.7 & 0.9 & 0 &  &  &  &  &  &  \\
5 & 4.4 & 1.3 & 0.7 & 1.1 & 0 &  &  &  &  &  \\
6 & 4.0 & 1.5 & 1.2 & 1.7 & 0.8 & 0 &  &  &  &  \\  
7 & 5.0 & 2.1 & 1.4 & 1.0 & 1.0 & 1.2 & 0 &  &  &  \\  
8 & 5.5 & 2.2 & 1.4 & 0.6 & 1.3 & 1.9 & 0.7 & 0 &  &  \\
9 & 5.7 & 3.0 & 2.2 & 1.4 & 2.2 & 2.6 & 1.2 & 0.9 & 0 &  \\   
10 & 5.2 & 3.3 & 2.7 & 2.2 & 2.6 & 2.9 & 1.9 & 1.7 & 1.1 & 0 \\          
\end{vmatrix}
 %
\hspace{3pt}
\begin{vmatrix}
    & 1 & 2 & 3 & 4 & 5 & 6 & 7 & 8 & 9 & 10 \\
1 & 0 &  &  &  &  &  &  &  &  &  \\
2 & 1.7 & 0 &  &  &  &  &  &  &  &  \\
3 & 3.3 & 1.1 & 0 &  &  &  &  &  &  &  \\
4 & 3.3 & 1.0 & 0.7 & 0 &  &  &  &  &  &  \\
5 & 4.2 & 1.9 & 0.9 & 1.2 & 0 &  &  &  &  &  \\
6 & 3.9 & 1.6 & 1.2 & 0.8 & 1.0 & 0 &  &  &  &  \\  
7 & 3.5 & 1.8 & 1.6 & 1.2 & 1.2 & 0.6 & 0 &  &  &  \\  
8 & 4.6 & 2.3 & 1.5 & 1.7 & 0.7 & 1.2 & 1.1 & 0 &  &  \\
9 & 5.0 & 3.0 & 2.2 & 2.5 & 1.5 & 1.9 & 1.6 & 0.9 & 0 &  \\   
10 & 3.5 & 2.6 & 2.2 & 2.3 & 1.8 & 2.0 & 1.7 & 1.4 & 0.9 & 0 \\              
\end{vmatrix} 
\end{equation}
\end{adjustwidth}
2: Bimonthly class separability matrices between 10 classes of land cover types in Rif Mountains, Morocco: May and July 2023 (2).

\begin{adjustwidth}{-\extralength}{0cm}
\begin{equation}
\scriptsize
\begin{vmatrix}
   & 1 & 2 & 3 & 4 & 5 & 6 & 7 & 8 & 9 & 10 \\
1 & 0 &  &  &  &  &  &  &  &  &  \\
2 & 2.8 & 0 &  &  &  &  &  &  &  &  \\
3 & 4.8 & 1.1 & 0 &  &  &  &  &  &  &  \\
4 & 3.6 & 1.0 & 0.6 & 0 &  &  &  &  &  &  \\
5 & 6.3 & 2.1 & 1.0 & 1.1 & 0 &  &  &  &  &  \\
6 & 6.8 & 2.3 & 1.3 & 1.1 & 0.6 & 0 &  &  &  &  \\  
7 & 8.2 & 3.1 & 2.1 & 1.7 & 1.1 & 0.8 & 0 &  &  &  \\  
8 & 8.6 & 3.6 & 2.6 & 2.1 & 1.8 & 1.4 & 0.7 & 0 &  &  \\
9 & 8.6 & 4.0 & 3.2 & 2.5 & 2.5 & 2.1 & 1.5 & 0.8 & 0 &  \\   
10 & 7.1 & 4.1 & 3.4 & 2.9 & 2.9 & 2.6 & 2.2 & 1.6 & 1.0 & 0 \\          
\end{vmatrix}
 %
\hspace{3pt}
\begin{vmatrix}
    & 1 & 2 & 3 & 4 & 5 & 6 & 7 & 8 & 9 & 10 \\
1 & 0 &  &  &  &  &  &  &  &  &  \\
2 & 2.0 & 0 &  &  &  &  &  &  &  &  \\
3 & 3.3 & 1.0 & 0 &  &  &  &  &  &  &  \\
4 & 3.9 & 1.6 & 0.9 & 0 &  &  &  &  &  &  \\
5 & 5.0 & 2.3 & 1.3 & 0.8 & 0 &  &  &  &  &  \\
6 & 3.6 & 1.6 & 0.8 & 1.0 & 0.9 & 0 &  &  &  &  \\  
7 & 5.5 & 2.8 & 1.8 & 1.6 & 0.8 & 1.0 & 0 &  &  &  \\  
8 & 6.0 & 3.2 & 2.2 & 1.8 & 1.0 & 1.5 & 0.7 & 0 &  &  \\
9 & 6.1 & 3.6 & 2.8 & 2.6 & 1.8 & 1.9 & 1.3 & 1.0 & 0 &  \\   
10 & 5.3 & 3.8 & 3.2 & 3.1 & 2.6 & 2.5 & 2.2 & 2.0 & 1.2 & 0 \\              
\end{vmatrix} 
\end{equation}
\end{adjustwidth}
3: Bimonthly class separability matrices between 10 classes of land cover types in Rif Mountains, Morocco: September and November 2023 (3).
\vspace{12pt}
%MDPI:  1. We have modified the format of the formula, please confirm%Author: yes, confirm / 2. Please confirm whether the nonindentation is maintained and whether need to be modified to normal indented paragraph. Same as the following text. %Author: paragraph can be with or without nonindentation (not so important). I agree with your correction.

\vspace{-8pt}
\begin{figure}[H]
    \begin{center}
	\hspace{-0pt}\includegraphics[width=13.7cm]{Fig_09}
    \end{center}
    \caption{Accuracy evaluation using computed correctly classified pixels (\%), GRASS GIS.}
    \label{fig09}
\end{figure}

The visualization of Earth landscapes using cartographic methods has many approaches; one of them is the use of RS data. This method uses satellite images as input data for their processing and classification. Typically, cartography relies on the use of traditional GIS methods that can produce a visualization using a defined workflow dependent on a particular software. The detection of land cover types depicted on these maps is perceptible to the human eye depending on the spatial and temporal resolution of the RS data. However, the use of automation and ML approaches facilitates image processes, thanks to a more objective and accurate classification of satellite images.

%----------- subsection ----------->
\subsection{{Bimonthly Evaluation of the Land Cover Types}}

{The bimonthly evaluation of land cover types offers an advantage over annual dynamics, where usually the landscape is evaluated at the same time but with intervals of several years. In contrast, bimonthly evaluations enable a more close look into the gradual slight changes that are caused by the seasonal climate effects. }Complex interactions between human activities, climate change, and environmental dynamics result in a high vulnerability of landscapes and ecosystems in Morocco. This is particularly important in northern Morocco, a region which is highly sensitive to the contrasting effects of climate variability, ranging from the arid Sahara to the humid Mediterranean, and where current land cover types are affected by pastoralism. As well as highlighting the current distribution of land cover types in northern Morocco, using RS data processed by ML algorithms presents new maps that allow for the quantification of land cover changes computed in pixels. For instance, numerical analysis has revealed that even for the least modified areas of bare areas located near desert, regional inventories show slight changes over diverse seasonal periods during the calendar year. Spatial biases are also apparent if comparing such land cover types as grassland, croplands, \mbox{and forests.}

\subsection{{Accuracy Assessment}}

{The spectral reflectance of these elements can be formally evaluated through digital numbers read from the satellite images using ML algorithms. Hence, selecting an optimal approach that accurately discriminates pixels on the image and identifies their category correctly presents an important process in cartographic framework. The suitability of algorithms has been assessed according to four criteria: (1) ease of constructing using functionality of GRASS GIS modules, (2) the ability of the algorithm to derive physical insight into land cover types from the satellite images for environmental investigation, (3) the ability to use each model to make quantitations based the processed satellite data and to generate numerical reports, (4) accuracy assessment in which the models might be tested.}

Here, the rejection threshold map contains the index of a confidence level calculated for each classified pixel in the classified image of the 10 land cover classes in the Morocco region. Confidence intervals are defined and the rejection map is visualized with lower values meaning 'keep the correctly classified pixel' and higher values meaning “reject a misclassified pixel that could belong to another class category”. Thus, this map identifies cells in the classified raster image that have a low probability, i.e., a high rejection index, of being correctly assigned to the target class.

The maps describing land use in the Rif mountain region are scientific and useful decision-making tools for Moroccan authorities and educational institutions. The images processed by the GRASS GIS methods can be used in research work as the input data for modeling systems of bimonthly vegetation cycles in the geomorphological context of the Rif. In addition, they can also be used for the operational monitoring of bimonthly vegetation changes in northern Morocco in order to apply research recommendations that require precise knowledge of the \mbox{Rif territories.}





\begin{table}[H]


%MDPI: Please cite the table in the text and ensure that the first citation of each figure appears in numerical order.
\caption{\hl{Bimonthly} evaluation of the class means of Digital Numbers (DN) computed for pixels assigned to 10 land cover classes in landscapes of Rif Mountains, north Morocco. Each multispectral band (1 to 7) of Landsat 8--9 OLI/TIRS is evaluated. Period: 2023, July to November.} \label{tab04}
\begin{tabularx}{\textwidth}{@{\extracolsep{\fill}}p{1.5cm}cccccccc@{}}
\toprule  \textbf{Class} & \textbf{Band 1} & \textbf{Band 2} & \textbf{Band 3} & \textbf{Band 4} & \textbf{Band 5} & \textbf{Band 6} & \textbf{Band 7} \\   \midrule

\multicolumn{8}{c }{\textbf{\hl{2023: July}}}\\
  \midrule
1 & 7500.81 & 7646.92 & 7722.75 & 7376.19 & 7464.38 & 7798.77 & 7755.02 \\
2 & 8127.96 & 8540.32 & 9609.1 & 9705.24 & 15,267.4 & 13,436.5 & 11,015.4 \\
3 & 8252.89 & 8754.25 & 10,217.3 & 10,220.6 & 19,793.6 & 15,491.7 & 12,155.2 \\
4 & 8632.51 & 9160.58 & 10,487.1 & 11,145.4 & 15,532 & 15,921.1 & 13,380.2 \\
5 & 8846.49 & 9479.29 & 11,046.4 & 11,830.1 & 17,492.5 & 17,638.3 & 14,481.8 \\
6 & 9201.59 & 9873.76 & 11,411 & 12,488 & 15,585.3 & 18,152.4 & 15,358.2 \\
7 & 9293.07 & 10,015.1 & 11,685.1 & 13,029.5 & 17,033.7 & 19,868.5 & 16,751.3 \\
8 & 10,080.1 & 10,993.2 & 12,891.2 & 14,295.8 & 17,321.4 & 20,518.4 & 17,159.1 \\
9 & 9985.18 & 10,842.3 & 12,815.2 & 14,613.6 & 18,222.5 & 22,332.9 & 19,188 \\
10 & 10,968.9 & 12,001.6 & 14,474.9 & 16,602.5 & 19,960 & 24,353.3 & 21,345.3 \\
  \midrule
%\bf{2022} &  &  &  & & & & \\
\multicolumn{8}{c }{\textbf{\hl{2023: September}}}\\
  \midrule
1 & 8057.65 & 8039.26 & 7748.32 & 7366.06 & 7277.01 & 7381.23 & 7381.57 \\
2 & 7726.42 & 7990.22 & 8786.58 & 8828.21 & 14,163.6 & 12,041.9 & 10,029.1 \\
3 & 8247.64 & 8609.77 & 9618.74 & 10,150.4 & 14,170.5 & 14,536.1 & 12,265.5 \\
4 & 8313.97 & 8631 & 9791.5 & 9795.11 & 18,703.7 & 14,820.3 & 11,736 \\
5 & 8759.38 & 9204.49 & 10,345.5 & 11,259.8 & 13,905.7 & 16,652 & 14,282 \\
6 & 8851.15 & 9357.15 & 10,724.4 & 11,653 & 16,282.7 & 17,375 & 14,521.9 \\
7 & 9289.65 & 9878.35 & 11,350.2 & 12,612.5 & 15,641.1 & 18,960.8 & 16,187 \\
8 & 9656.35 & 10,336.8 & 12,032.2 & 13,588.6 & 16,911.5 & 20,589.9 & 17,615.2 \\
9 & 10,122.7 & 10,879.8 & 12,859.4 & 14,816 & 18,280.5 & 22,704 & 19,758.4 \\
10 & 11,172.9 & 12,165.6 & 14,822.7 & 17,351.5 & 20,822.6 & 25,700.6 & 22,812.4 \\
  \midrule

%\bf{2023} &  &  &  & & & & \\
\multicolumn{8}{c }{\textbf{\hl{2023: November}}}\\
  \midrule
1 & 7982.57 & 7976.11 & 7746.84 & 7371.42 & 7374.01 & 7416.82 & 7387.47 \\
2 & 7606.96 & 7848.7 & 8421.42 & 8450.72 & 11,883.1 & 10,576.5 & 9345.26 \\
3 & 8009.92 & 8271.76 & 9076.5 & 9113.51 & 14,879.4 & 13,007.5 & 10,800.5 \\
4 & 8443.29 & 8714.75 & 9430.96 & 9844.52 & 11,725.6 & 13,884.5 & 12,143.8 \\
5 & 8688.21 & 9043.79 & 10003.6 & 10602.4 & 13,637.5 & 15,634.6 & 13,570.2 \\
6 & 8371.09 & 8671.7 & 9834.34 & 9692.7 & 18,571.4 & 14,782.6 & 11,792.6 \\
7 & 8804.19 & 9223.59 & 10,496.4 & 11,154 & 16,783.6 & 17,135.9 & 14,444 \\
8 & 9160.99 & 9627.78 & 10,860.9 & 11,853.9 & 14,513.2 & 17,931.3 & 15,848 \\
9 & 9542.27 & 10,134.7 & 11,817.5 & 13,258.5 & 17,309 & 20,476.2 & 18,053.6 \\
10 & 11,042.5 & 12,012.4 & 14,554.4 & 16,720.1 & 20,561.1 & 25,084.9 & 22,211.2 \\
 \bottomrule
\end{tabularx}
\end{table}
%MDPI: Please confirm if the bold is unnecessary and can be removed. The following highlights are the same.

%----------- subsection ----------->



\vspace{-5pt}


\begin{table}[H]
\caption{\hl{Computed} points stability for iterations during classification and convergence of classes: January, March, May, July, September and November 2023.}\label{tabA06}
\begin{tabularx}{\textwidth}{p{1.7cm}  XX  X  X  X X }
  \toprule
   \textbf{Iteration} & \textbf{01/2023} & \textbf{03/2023} & \textbf{05/2023} & \textbf{07/2023} & \textbf{09/2023} & \textbf{11/2023}  \\
   \midrule
    Iteration 1 & 74.15\% & 81.46\% & 58.44\% & 74.15\% & 81.12\% & 81.47\% \\
    Iteration 2 & 77.52\% & 88.60\% & 74.76\% & 77.52\% & 88.74\% & 78.51\% \\
    Iteration 3 & 90.02\% & 88.06\% & 89.85\% & 90.02\% & 92.05\% & 85.56\% \\
    Iteration 4 & 93.59\% & 93.42\% & 94.81\% & 93.59\% & 94.67\% & 92.87\% \\
    Iteration 5 & 95.31\% & 95.40\% & 96.52\% & 95.31\% & 95.25\% & 96.37\% \\
    Iteration 6 & 96.32\% & 95.94\% & 96.85\% & 96.32\% & 96.21\% & 97.33\% \\
    Iteration 7 & 97.11\% & 96.88\% & 97.34\% & 97.11\% & 97.09\% & 97.66\% \\
    Iteration 8 & 97.05\% & 97.78\% & 97.60\% & 97.05\% & 97.19\% & 97.97\% \\
    Iteration 9 & 97.55\% & 98.16\% & 98.03\% & 97.55\% & 97.46\% & 98.24\% \\
    \midrule
    Convergence & 98.1\% & 98.2\% & 98.0\% & 98.1\% & 98.0\% & 98.2\% \\
    \bottomrule
\end{tabularx}

\end{table}
%MDPI: Please cite the table in the text and ensure that the first citation of each figure appears in numerical order.
%----------- section ----------->
\section{Conclusions}

The general objective of this paper was to improve the production of land cover maps in northern Morocco using ML methods of GRASS GIS. Our results clearly show that the land cover types in northern Morocco are changing. Nevertheless, the dynamics of such changes are vastly underrepresented in existing case studies on environmental mapping in northern Africa. At the same time, satellite images present the largest provider of information on the distribution of landscape patches and vegetation. To this end, we applied the advanced algorithms `Decision trees' and `Extra trees' for processing a time series of the Landsat 8-9 OLI/TIRS satellite images. 

More than this, however, we have also shown the technical difference between the unsupervised and supervised classification (clustering) of RS data. This is useful for the further development of cartographic methods of RS data processing in the paradigm of information processing systems. Therefore, the specific objective of this paper aims to improve land cover mapping in northern Morocco from the time series of Landsat satellite images using ML methods by GRASS GIS and cartographic scripts. Two challenges have been identified in this work. The first concerns the technical choice of the classification algorithm (supervised “DecisionTreeClassifier” and “ExtraTreesClassifier” vs unsupervised clustering by MaxLike) using GRASS GIS, while the second concerns the identification of seasonal bimonthly landscape dynamics in the Rif Mountains of Morocco using repetitive data cycles from January to November 2023. A comparison of these data enabled the relationship between the growth of vegetation and mountain geomorphology to be revealed.

Taken together, this paper shows that the least well recorded region of the northern Morocco is largely affected by environmental changes: the comparison of presented maps showed the dynamics of agricultural and forest areas in the order of 7--10\% within the year of 2023, and a slight increase in bare soil of 8\% during the same year. Moreover, a notable change in the areas occupied by deciduous and mixed forest is detected. Currently, there are no land cover maps available for Morocco that depict bimonthly dynamics of land cover types for the region of the Rif Mountains. In this regard, this paper contributes to improve this gap through presenting new maps for further evaluating of biophysical and environmental parameters in Northern Morocco.

This paper demonstrated the implementation of the ML cartographic approach to RS data processing, technically performed using GRASS GIS tools. The presented work enabled the building of a classification workflow of six multispectral satellite images covering Northern Morocco, around the Rif Mountains. Several applications of the presented results and the GRASS GIS satellite image processing methods are described in this paper. The demonstrated maps show the dynamics in the land cover of northern Morocco, which presents the environmental dynamics that reflect the impact factors from regional ecological, vegetal, and geological settings in North Africa.

%----------- section ----------->
%\authorcontributions{Supervision, conceptualization, methodology, software, resources, funding acquisition, and project administration, writing---original draft preparation, methodology, software, data curation, visualization, formal analysis, validation, writing---review and editing, and investigation, P.L.}

\funding{The publication was funded by the Editorial Office of Geomatics, Multidisciplinary Digital Publishing Institute (MDPI), who provided a 100\% discount for the APC of this manuscript.}

\institutionalreview{Not applicable.}

\informedconsent{Not applicable.}

\dataavailability{\hl{Not applicable}.} 
%MDPI: "Not applicable" is only used for review papers or articles for which no new data were created. Normally, the DAS for research articles can state "Data are contained within the article." or "Data are contained within the article and supplementary materials." Please refer to the complete guideline at https://www.mdpi.com/ethics#_bookmark21.


\acknowledgments{The authors thank the reviewers for reading and reviewing this manuscript.}

\conflictsofinterest{The author declares no conflicts of interest.} 

\vspace{444pt}

\abbreviations{Abbreviations}{
The following abbreviations are used in this manuscript:\\

\noindent 
\begin{tabular}{@{}ll}
API & Application Programming Interface \\
DTC & Decision Tree Classifier \\
ETC & Extra Trees Classifier \\
EO & Earth Observation \\
FAO & Food and Agriculture Organization \\
GRASS & Geographic Resources Analysis Support System \\
GIS & Geographic Information System \\
Landsat OLI/TIRS & Landsat Operational Land Imager and Thermal Infrared Sensor  \\
GEBCO & General Bathymetric Chart of the Oceans \\
GMT & Generic Mapping Tools \\
QGIS & Quantum Geographic Information System \\
ML & Machine Learning \\
RS & Remote Sensing \\
SWIR & Shortwave Infrared \\
NIR & Near Infrared \\
USGS & United States Geological Survey  \\
WGS84 & World Geodetic System 84 \\
\end{tabular}
}

%%%%%%%%%%%%%%%%%%%%%%%%%%%%%%%%%%%%%%%%%%
%% Optional
\appendixtitles{yes} % Leave argument "no" if all appendix headings stay EMPTY (then no dot is printed after "Appendix A"). If the appendix sections contain a heading then change the argument to "yes".
\appendixstart
\appendix
%--------------------------->
\section[\appendixname~\thesection]{\hl{Initial} Means for Classes of Pixels in Each Multispectral Band of the Landsat OLI/TIRS Satellite Images}
%MDPI: Appendixes should be cited. An appendix’s Figures/Tables/Equations can be cited in the main text, but are not mandatory, please check and revise

\vspace{-3pt}
\begin{table}[H]
\caption{January 2023: initial means for each multispectral band of Landsat OLI/TIRS images.}

\begin{tabularx}{\textwidth}{ p{1.5cm} X  X  X X  X  XX }
\toprule \textbf{Class} & \textbf{Band 1} & \textbf{Band 2} & \textbf{Band 3} & \textbf{Band 4} & \textbf{Band 5} & \textbf{Band 6} & \textbf{Band 7} \\  \midrule
    Class 1 & 7353.24 & 7533.21 & 8177.38 & 7818.66 & 10987.1 & 10158.3 & 8739.9 \\
    Class 2 & 7569.99 & 7774.52 & 8500.84 & 8237.08 & 11888.8 & 10959.1 & 9396.37 \\
    Class 3 & 7786.75 & 8015.83 & 8824.3 & 8655.49 & 12790.5 & 11759.9 & 10052.8 \\
    Class 4 & 8003.51 & 8257.14 & 9147.76 & 9073.91 & 13692.2 & 12560.7 & 10709.3 \\
    Class 5 & 8220.27 & 8498.45 & 9471.21 & 9492.32 & 14593.9 & 13361.5 & 11365.8 \\
    Class 6 & 8437.03 & 8739.76 & 9794.67 & 9910.73 & 15495.6 & 14162.3 & 12022.2 \\
    Class 7 & 8653.78 & 8981.07 & 10,118.1 & 10329.1 & 16,397.3 & 14,963.1 & 12,678.7 \\
    Class 8 & 8870.54 & 9222.38 & 10,441.6 & 10,747.6 & 17,299 & 15,763.9 & 13,335.2 \\
    Class 9 & 9087.3 & 9463.69 & 10765 & 11,166 & 18,200.8 & 16,564.7 & 13,991.6 \\
    Class 10 & 9304.06 & 9705 & 11,088.5 & 11,584.4 & 19,102.5 & 17,365.5 & 14,648.1 \\
    \bottomrule
\end{tabularx}
\label{tabA01}
\end{table}


\vspace{-6pt}
\begin{table}[H]
\caption{March 2023: initial means for each multispectral band of Landsat OLI/TIRS satellite images.}\label{tabA02}
\begin{tabularx}{\textwidth}{ p{1.5cm}  X  X  X  X  X  X  X  }
\toprule \textbf{Class} & \textbf{Band 1} & \textbf{Band 2} & \textbf{Band 3} & \textbf{Band 4} & \textbf{Band 5} & \textbf{Band 6} & \textbf{Band 7} \\  \midrule
    Class 1 & 8036.12 & 8397.21 & 9295.02 & 9561.14 & 13138.1 & 12962.9 & 10897.1 \\
    Class 2 & 8246.35 & 8665.51 & 9702.28 & 10,172.8 & 14,019.7 & 13,972.2 & 11,692.2 \\
    Class 3 & 8456.57 & 8933.81 & 10,109.5 & 10,784.5 & 14,901.3 & 14,981.5 & 12,487.3 \\
    Class 4 & 8666.8 & 9202.11 & 10,516.8 & 11,396.2 & 15,783 & 15,990.8 & 13,282.5 \\
    Class 5 & 8877.03 & 9470.41 & 10,924.1 & 12,007.9 & 16,664.6 & 17,000.2 & 14,077.6 \\
    Class 6 & 9087.26 & 9738.71 & 11,331.3 & 12,619.6 & 17,546.2 & 18,009.5 & 14,872.7 \\
    Class 7 & 9297.49 & 10,007 & 11,738.6 & 13,231.3 & 18,427.9 & 19,018.8 & 15,667.8 \\
    Class 8 & 9507.72 & 10,275.3 & 12,145.8 & 13,842.9 & 19,309.5 & 20,028.1 & 16,463 \\
    Class 9 & 9717.94 & 10,543.6 & 12,553.1 & 14,454.6 & 20,191.1 & 21,037.4 & 17,258.1 \\
    Class 10 & 9928.17 & 10,811.9 & 12,960.4 & 15,066.3 & 21,072.7 & 22,046.7 & 18,053.2 \\
    \bottomrule
\end{tabularx}
\label{tabA02}
\end{table}

\begin{table}[H]
\caption{May 2023: initial means for each multispectral band of Landsat OLI/TIRS satellite images.}\label{tabA03}
\begin{tabularx}{\textwidth}{ p{1.5cm}  X  X  X  X  X  X  X  }
\toprule \textbf{Class} & \textbf{Band 1} & \textbf{Band 2} & \textbf{Band 3} & \textbf{Band 4} & \textbf{Band 5} & \textbf{Band 6} & \textbf{Band 7} \\  \midrule
    Class 1 & 7598.78 & 7764.07 & 8434.53 & 7941.99 & 12,295.6 & 10,618.7 & 8849.61 \\
    Class 2 & 7754.07 & 7949.13 & 8718.98 & 8335.56 & 13,270.4 & 11,401 & 9511.46 \\
    Class 3 & 7909.36 & 8134.19 & 9003.43 & 8729.14 & 14,245.2 & 12,183.3 & 10,173.3 \\
    Class 4 & 8064.65 & 8319.25 & 9287.89 & 9122.72 & 15,219.9 & 12,965.7 & 10,835.2 \\
    Class 5 & 8219.95 & 8504.31 & 9572.34 & 9516.29 & 16,194.7 & 13,748 & 11,497 \\
    Class 6 & 8375.24 & 8689.37 & 9856.79 & 9909.87 & 17,169.5 & 14,530.3 & 12,158.8 \\
    Class 7 & 8530.53 & 8874.43 & 10,141.2 & 10,303.4 & 18,144.3 & 15,312.6 & 12,820.7 \\
    Class 8 & 8685.82 & 9059.49 & 10,425.7 & 10,697 & 19,119.1 & 16,095 & 13,482.5 \\
    Class 9 & 8841.12 & 9244.55 & 10,710.1 & 11,090.6 & 20,093.9 & 16,877.3 & 14,144.4 \\
    Class 10 & 8996.41 & 9429.61 & 10,994.6 & 11,484.2 & 21,068.7 & 17,659.6 & 14,806.2 \\
    \bottomrule
\end{tabularx}
\label{tabA03}
\end{table}
\vspace{-6pt}


\begin{table}[H]
\caption{July 2023: initial means for each multispectral band of Landsat OLI/TIRS satellite images.}\label{tabA04}
\begin{tabularx}{\textwidth}{ p{1.5cm} X X X X X X X }
\toprule \textbf{Class} & \textbf{Band 1} & \textbf{Band 2} & \textbf{Band 3} & \textbf{Band 4} & \textbf{Band 5} & \textbf{Band 6} & \textbf{Band 7} \\  \midrule
    Class 1 & 8146.96 & 8572.74 & 9528.06 & 9871.64 & 12,802.2 & 13,459.3 & 11,517.8 \\
    Class 2 & 8370.6 & 8848.08 & 9926.18 & 10,412.8 & 13,530.9 & 14,431.9 & 12,322.4 \\
    Class 3 & 8594.23 & 9123.41 & 10,324.3 & 10,954 & 14,259.6 & 15,404.5 & 13,127.1 \\
    Class 4 & 8817.87 & 9398.75 & 10,722.4 & 11,495.2 & 14,988.3 & 16,377.2 & 13,931.7 \\
    Class 5 & 9041.5 & 9674.09 & 11,120.5 & 12,036.4 & 15,717.1 & 17,349.8 & 14,736.3 \\
    Class 6 & 9265.14 & 9949.43 & 11,518.6 & 12,577.6 & 16,445.8 & 18,322.4 & 15,541 \\
    Class 7 & 9488.77 & 10,224.8 & 11,916.7 & 13,118.9 & 17,174.5 & 19,295.1 & 16,345.6 \\
    Class 8 & 9712.41 & 10,500.1 & 12,314.9 & 13,660.1 & 17,903.2 & 20,267.7 & 17,150.3 \\
    Class 9 & 9936.04 & 10,775.4 & 12,713 & 14,201.3 & 18,631.9 & 21,240.3 & 17,954.9 \\
    Class 10 & 10,159.7 & 11,050.8 & 13,111.1 & 14,742.5 & 19,360.6 & 22,212.9 & 18,759.6 \\
    \bottomrule
\end{tabularx}
\label{tabA04}
\end{table}
\vspace{-6pt}
\begin{table}[H]
\caption{September 2023: initial means for each multispectral band of Landsat OLI/TIRS images.}\label{tabA05}
\begin{tabularx}{\textwidth}{ p{1.5cm} X X X X X X X }
\toprule \textbf{Class} & \textbf{Band 1} & \textbf{Band 2} & \textbf{Band 3} & \textbf{Band 4} & \textbf{Band 5} & \textbf{Band 6} & \textbf{Band 7} \\  \midrule
    Class 1 & 8037.45 & 8299.79 & 8964.58 & 9252.55 & 11,892.9 & 12,555.5 & 10,819.7 \\
    Class 2 & 8252.53 & 8565.39 & 9363.68 & 9802.02 & 12,631.3 & 13,569.5 & 11,680.2 \\
    Class 3 & 8467.62 & 8830.99 & 9762.77 & 10,351.5 & 13,369.8 & 14,583.6 & 12,540.8 \\
    Class 4 & 8682.7 & 9096.59 & 10,161.9 & 10,901 & 14,108.2 & 15,597.6 & 13,401.3 \\
    Class 5 & 8897.79 & 9362.19 & 10,561 & 11,450.4 & 14,846.6 & 16,611.6 & 14,261.9 \\
    Class 6 & 9112.87 & 9627.78 & 10,960 & 11,999.9 & 15,585.1 & 17,625.6 & 15,122.4 \\
    Class 7 & 9327.96 & 9893.38 & 11,359.1 & 12,549.4 & 16,323.5 & 18,639.6 & 15,983 \\
    Class 8 & 9543.04 & 10,159 & 11,758.2 & 13,098.9 & 17,061.9 & 19,653.6 & 16,843.5 \\
    Class 9 & 9758.13 & 10,424.6 & 12,157.3 & 13,648.3 & 17,800.4 & 20,667.6 & 17,704.1 \\
    Class 10 & 9973.21 & 10,690.2 & 12,556.4 & 14,197.8 & 18,538.8 & 21,681.7 & 18,564.6 \\
    \bottomrule
\end{tabularx}
\label{tabA05}
\end{table}
\vspace{-6pt}
\begin{table}[H]
\caption{November 2023: initial means for each multispectral band of Landsat OLI/TIRS images.}\label{tabA06}
\begin{tabularx}{\textwidth}{p{1.5cm} X X X X X X X }
\toprule \textbf{Class} & \textbf{Band 1} & \textbf{Band 2} & \textbf{Band 3} & \textbf{Band 4} & \textbf{Band 5} & \textbf{Band 6} & \textbf{Band 7} \\  \midrule
    Class 1 & 7805.3 & 7978.69 & 8394.85 & 8339.26 & 10,618.8 & 10,876.8 & 9549.69 \\
    Class 2 & 7994.42 & 8207.25 & 8745.73 & 8818 & 11,398.4 & 11,816.7 & 10,351.8 \\
    Class 3 & 8183.55 & 8435.81 & 9096.6 & 9296.74 & 12,178 & 12,756.6 & 11,154 \\
    Class 4 & 8372.67 & 8664.38 & 9447.48 & 9775.48 & 12,957.6 & 13,696.6 & 11,956.1 \\
    Class 5 & 8561.79 & 8892.94 & 9798.36 & 10,254.2 & 13,737.2 & 14,636.5 & 12,758.2 \\
    Class 6 & 8750.92 & 9121.5 & 10,149.2 & 10,733 & 14,516.8 & 15,576.5 & 13,560.4 \\
    Class 7 & 8940.04 & 9350.07 & 10,500.1 & 11,211.7 & 15,296.4 & 16,516.4 & 14,362.5 \\
    Class 8 & 9129.17 & 9578.63 & 10,851 & 11,690.4 & 16,076 & 17,456.4 & 15,164.7 \\
    Class 9 & 9318.29 & 9807.2 & 11,201.9 & 12,169.2 & 16,855.6 & 18,396.3 & 15,966.8 \\
    Class 10 & 9507.41 & 10,035.8 & 11,552.8 & 12,647.9 & 17,635.2 & 19,336.3 & 16,768.9 \\
    \bottomrule
\end{tabularx}
\label{tabA06}
\end{table}

%--------------------------->
\section[\appendixname~\thesection]{Standard Deviations for DNs of Pixels by 10 Computed Classes}
\vspace{-3pt}
\begin{table}[H]
\caption{Bimonthly evaluation of the standard deviations of Digital Numbers (DN) computed for pixels assigned to 10 land cover classes in landscapes of Rif Mountains, north Morocco. Evaluated each multispectral band (1 to 7) of Landsat 8--9 OLI/TIRS. Period: 2023, January to May.} \label{tabA2}

\begin{tabularx}{\textwidth}{@{\extracolsep{\fill}}p{1.5cm}cccccccc@{}}

\toprule \textbf{Class} & \textbf{Band 1} & \textbf{Band 2} & \textbf{Band 3} & \textbf{Band 4} & \textbf{Band 5} & \textbf{Band 6} & \textbf{Band 7} \\  \midrule




%\bf{2019} &  &  &  &  &  &  & \\
\multicolumn{8}{ c }{\textbf{2023: January}}\\
 \midrule
1 & 520.366 & 535.906 & 762.029 & 534.589 & 458.971 & 261.960 & 185.243 \\
2 & 490.239 & 471.037 & 624.903 & 718.469 & 1108.75 & 864.994 & 643.283 \\
3 & 469.625 & 475.629 & 523.194 & 551.304 & 924.17 & 813.025 & 684.606 \\
4 & 314.352 & 297.355 & 402.396 & 473.559 & 946.247 & 881.536 & 697.401 \\
5 & 355.762 & 363.823 & 482.001 & 628.670 & 1123.910 & 809.348 & 943.121 \\
6 & 398.402 & 397.373 & 474.544 & 609.048 & 1059.400 & 1019.420 & 849.959 \\
7 & 831.164 & 913.648 & 997.848 & 1243.020 & 1609.220 & 1416.340 & 1089.100 \\
8 & 602.747 & 649.726 & 694.899 & 843.584 & 1472.750 & 926.428 & 972.410 \\
9 & 563.373 & 645.449 & 846.794 & 1005.060 & 1844.240 & 1204.890 & 1309.140 \\
10 & 3912.350 & 4022.450 & 3523.310 & 3432.850 & 3098.910 & 2796.470 & 2599.350 \\
 \midrule
%\bf{2022} &  &  &  & & & & \\
\multicolumn{8}{ c }{\textbf{2023: March}}\\
 \midrule
1 & 490.469 & 541.212 & 879.207 & 716.637 & 616.488 & 686.935 & 613.483 \\
2 & 255.312 & 246.682 & 318.072 & 462.128 & 1664.020 & 929.262 & 719.619 \\
3 & 541.560 & 602.780 & 686.249 & 782.835 & 1273.070 & 980.419 & 834.318 \\
4 & 344.733 & 387.283 & 524.843 & 812.884 & 1659.040 & 980.065 & 855.260 \\
5 & 377.657 & 422.529 & 565.366 & 724.443 & 1044.500 & 753.815 & 856.870 \\
6 & 507.082 & 559.630 & 652.069 & 974.570 & 1023.500 & 866.568 & 784.542 \\
7 & 449.736 & 488.873 & 597.974 & 692.265 & 1060.560 & 878.137 & 929.794 \\
8 & 532.728 & 589.376 & 727.215 & 948.085 & 1317.450 & 934.689 & 1002.260 \\
9 & 507.096 & 562.500 & 709.818 & 853.806 & 1051.740 & 1083.610 & 1336.400 \\
10 & 753.846 & 857.879 & 1039.350 & 1108.080 & 1502.720 & 1502.480 & 2189.920 \\
 \midrule

%\bf{2023} &  &  &  & & & & \\
\multicolumn{8}{ c }{\textbf{2023: May}}\\
 \midrule
1 & 465.095 & 566.529 & 901.861 & 522.758 & 321.154 & 193.603 & 120.269 \\
2 & 277.062 & 255.083 & 376.608 & 450.477 & 1258.88 & 804.346 & 710.667 \\
3 & 441.646 & 468.557 & 562.157 & 597.613 & 1136.880 & 723.626 & 673.917 \\
4 & 378.651 & 385.161 & 519.410 & 649.726 & 1220.780 & 868.215 & 1039.880 \\
5 & 258.976 & 264.420 & 393.296 & 494.295 & 1055.540 & 951.268 & 799.838 \\
6 & 228.533 & 221.204 & 444.333 & 414.391 & 1722.940 & 994.491 & 630.409 \\
7 & 531.397 & 586.941 & 678.976 & 806.428 & 1337.73 & 1061.7 & 936.166 \\
8 & 623.942 & 671.415 & 760.778 & 817.799 & 1089.700 & 905.306 & 884.892 \\
9 & 528.555 & 602.147 & 767.459 & 860.058 & 1806.600 & 1088.830 & 1165.730 \\
10 & 1210.190 & 1398.230 & 1656.380 & 1682.400 & 1778.850 & 1867.540 & 2167.590 \\
\bottomrule
\end{tabularx}
\end{table}

\begin{table}[H]
\caption{Bimonthly evaluation of the standard deviations of Digital Numbers (DN) computed for pixels assigned to 10 land cover classes in landscapes of Rif Mountains, north Morocco. Evaluated each multispectral band (1 to 7) of Landsat 8--9 OLI/TIRS. Period: 2023, July to November.} \label{tab04}


\begin{tabularx}{\textwidth}{@{\extracolsep{\fill}}p{1.5cm}cccccccc@{}}

\toprule \textbf{Class} & \textbf{Band 1} & \textbf{Band 2} & \textbf{Band 3} & \textbf{Band 4} & \textbf{Band 5} & \textbf{Band 6} & \textbf{Band 7} \\  \midrule 


%\bf{2019} &  &  &  &  &  &  & \\
\multicolumn{8}{ c }{\textbf{2023: July}}\\
 \midrule 
1 & 401.603 & 474.680 & 753.823 & 554.671 & 452.066 & 377.558 & 279.117 \\
2 & 444.544 & 494.063 & 605.472 & 689.234 & 1469.020 & 1074.000 & 858.682 \\
3 & 400.815 & 449.109 & 548.777 & 721.564 & 1773.570 & 1030.710 & 912.245 \\
4 & 329.822 & 337.588 & 456.271 & 617.552 & 934.712 & 737.397 & 693.313 \\
5 & 383.645 & 392.394 & 484.058 & 558.823 & 955.127 & 757.462 & 640.994 \\
6 & 369.823 & 422.794 & 574.353 & 625.855 & 807.736 & 699.93 & 634.601 \\
7 & 380.957 & 393.860 & 462.246 & 493.918 & 938.201 & 756.140 & 818.874 \\
8 & 389.170 & 433.991 & 512.028 & 543.331 & 810.822 & 715.835 & 672.958 \\
9 & 585.918 & 680.160 & 803.684 & 738.175 & 743.199 & 762.175 & 944.871 \\
10 & 994.099 & 1126.100 & 1284.760 & 1153.090 & 1115.620 & 1242.490 & 1575.490 \\
 \midrule 
%\bf{2022} &  &  &  & & & & \\
\multicolumn{8}{ c }{\textbf{2023: September}}\\
 \midrule 
1 & 437.651 & 443.108 & 680.601 & 464.861 & 251.267 & 134.622 & 109.822 \\
2 & 352.790 & 343.761 & 458.110 & 604.993 & 1647.410 & 1009.150 & 737.153 \\
3 & 413.844 & 428.298 & 527.591 & 627.276 & 1270.460 & 820.868 & 669.725 \\
4 & 609.427 & 665.603 & 750.628 & 922.569 & 2049.290 & 1237.420 & 1107.210 \\
5 & 410.913 & 444.970 & 557.272 & 605.301 & 936.943 & 802.527 & 790.452 \\
6 & 412.096 & 463.575 & 605.626 & 711.301 & 1026.750 & 826.580 & 710.861 \\
7 & 476.481 & 556.664 & 690.121 & 675.737 & 961.119 & 727.211 & 750.357 \\
8 & 560.709 & 661.131 & 788.185 & 748.966 & 878.300 & 826.902 & 929.749 \\
9 & 734.715 & 872.060 & 1021.830 & 874.668 & 875.353 & 1032.820 & 1289.020 \\
10 & 1228.510 & 1477.600 & 1759.470 & 1582.050 & 1479.630 & 1594.770 & 1825.860 \\
 \midrule 

%\bf{2023} &  &  &  & & & & \\
\multicolumn{8}{ c}{\textbf{2023: November}}\\
 \midrule 
1 & 408.799 & 366.029 & 518.468 & 358.299 & 494.063 & 288.290 & 183.250 \\
2 & 337.779 & 313.223 & 378.857 & 471.889 & 1236.690 & 864.664 & 718.244 \\
3 & 329.228 & 309.538 & 391.118 & 526.252 & 1133.950 & 971.402 & 842.809 \\
4 & 357.573 & 340.276 & 405.652 & 474.368 & 868.293 & 878.208 & 719.913 \\
5 & 407.823 & 418.715 & 501.585 & 583.480 & 1177.420 & 746.789 & 760.614 \\
6 & 560.647 & 596.644 & 680.376 & 807.904 & 1923.630 & 1063.100 & 992.282 \\
7 & 404.268 & 434.357 & 567.739 & 689.356 & 1145.220 & 911.636 & 828.516 \\
8 & 502.947 & 558.014 & 686.296 & 677.680 & 976.150 & 885.420 & 857.641 \\
9 & 571.593 & 667.463 & 845.177 & 950.398 & 1288.210 & 1164.560 & 1261.840 \\
10 & 1130.670 & 1394.490 & 1762.720 & 1842.670 & 1953.210 & 2219.250 & 2058.000 \\
\bottomrule
\end{tabularx}
\end{table}

%%%%%%%%%%%%%%%%%%%%%%%%%%%%%%%%%%%%%%%%%%
\begin{adjustwidth}{-\extralength}{0cm}
%\printendnotes[custom] % Un-comment to print a list of endnotes

\reftitle{References}
%\nocite{*}

%=====================================
% References, variant A: external bibliography
%=====================================
\begin{thebibliography}{999}

\bibitem[Corenblit and Steiger(2009)]{Corenblit}
Corenblit, D.; Steiger, J.
\newblock Vegetation as a major conductor of geomorphic changes on the Earth
  surface: Toward evolutionary geomorphology.
\newblock {\em Earth Surf. Process. Landforms} {\bf 2009}, {\em
  34},~891--896.
\newblock {\url{https://doi.org/10.1002/esp.1788}}.

\bibitem[Muir et~al.(2024)Muir, Hurst, Richardson-Foulger, Rennie, and
  Naylor]{Muir}
Muir, F.M.E.; Hurst, M.D.; Richardson-Foulger, L.; Rennie, A.F.; Naylor, L.A.
\newblock VedgeSat: An automated, open-source toolkit for coastal change
  monitoring using satellite-derived vegetation edges.
\newblock {\em Earth Surf. Process. Landforms} {\bf 2024}, {\em
  49},~2405--2423.
\newblock {\url{https://doi.org/10.1002/esp.5835}}.

\bibitem[Pu et~al.(2021)Pu, Quackenbush, and Stehman]{Pu}
Pu, G.; Quackenbush, L.J.; Stehman, S.V.
\newblock Using Google Earth Engine to Assess Temporal and Spatial Changes in
  River Geomorphology and Riparian Vegetation.
\newblock {\em JAWRA J. Am. Water Resour. Assoc.} {\bf
  2021}, {\em 57},~789--806.
\newblock {\url{https://doi.org/10.1111/1752-1688.12950}}.

\bibitem[Cary(1990)]{Cary01091990}
Cary, T.
\newblock Landsat Data as a Tool for the Geosciences.
\newblock {\em J. Geol. Educ.} {\bf 1990}, {\em 38},~318--322.
\newblock {\url{https://doi.org/10.5408/0022-1368-38.4.318}}.

\bibitem[Xiong et~al.(2023)Xiong, Chen, Li, Su, Dai, and Wang]{rs15225374}
Xiong, N.; Chen, H.; Li, R.; Su, H.; Dai, S.; Wang, J.
\newblock A Method of Chestnut Forest Identification Based on Time Series and
  Key Phenology from Sentinel-2.
\newblock {\em Remote. Sens.} {\bf 2023}, {\em 15}, 5374.
\newblock {\url{https://doi.org/10.3390/rs15225374}}.

\bibitem[Chen et~al.(2025)Chen, Zhang, Zhuang, and Hu]{w17010068}
Chen, X.; Zhang, X.; Zhuang, C.; Hu, X.
\newblock Detecting the Lake Area Seasonal Variations in the Tibetan Plateau
  from Multi-Sensor Satellite Data Using Deep Learning.
\newblock {\em Water} {\bf 2025}, {\em 17}, 68.
\newblock {\url{https://doi.org/10.3390/w17010068}}.

\bibitem[Kuta et~al.(2025)Kuta, Grebby, and Boyd]{app15010338}
Kuta, A.A.; Grebby, S.; Boyd, D.S.
\newblock Remote Monitoring of the Impact of Oil Spills on Vegetation in the
  Niger Delta, Nigeria.
\newblock {\em Appl. Sci.} {\bf 2025}, {\em 15}, 338.
\newblock {\url{https://doi.org/10.3390/app15010338}}.

\bibitem[Pan et~al.(2023)Pan, Li, Yang, Ren, Ma, Chen, Feng, Chen, Chen, and
  Li]{rs15225413}
Pan, D.; Li, C.; Yang, G.; Ren, P.; Ma, Y.; Chen, W.; Feng, H.; Chen, R.; Chen,
  X.; Li, H.
\newblock Identification of the Initial Anthesis of Soybean Varieties Based on
  UAV Multispectral Time-Series Images.
\newblock {\em Remote. Sens.} {\bf 2023}, {\em 15}, 5413.
\newblock {\url{https://doi.org/10.3390/rs15225413}}.

\bibitem[Ding and Li(2023)]{rs15225408}
Ding, N.; Li, M.
\newblock Mapping Forest Abrupt Disturbance Events in Southeastern
  China---Comparisons and Tradeoffs of Landsat Time Series Analysis Algorithms.
\newblock {\em Remote. Sens.} {\bf 2023}, {\em 15}, 5408.
\newblock {\url{https://doi.org/10.3390/rs15225408}}.

\bibitem[Yang et~al.(2023)Yang, Hu, Ma, Wang, Zhang, Li, and Zhang]{rs15225415}
Yang, J.; Hu, Y.; Ma, Y.; Wang, M.; Zhang, N.; Li, Z.; Zhang, J.
\newblock Combined Retrieval of Oil Film Thickness Using Hyperspectral and
  Thermal Infrared Remote Sensing.
\newblock {\em Remote. Sens.} {\bf 2023}, {\em 15}, 5415.
\newblock {\url{https://doi.org/10.3390/rs15225415}}.

\bibitem[Nima~Ahmadian and Zölitz(2016)]{Ahmadian03052016}
Nima~Ahmadian, Sahar~Ghasemi, J.P.W.; Zölitz, R.
\newblock Comprehensive study of the biophysical parameters of agricultural
  crops based on assessing Landsat 8 OLI and Landsat 7 ETM+ vegetation indices.
\newblock {\em GIsci. Remote. Sens.} {\bf 2016}, {\em 53},~337--359.
\newblock {\url{https://doi.org/10.1080/15481603.2016.1155789}}.

\bibitem[Lemenkova(2024)]{Lemenkova202412}
Lemenkova, P.
\newblock {Mapping Woodlands in Angola, Tropical Africa: Calculation of
  Vegetation Indices From Remote Sensing Data}.
\newblock {\em Agric. For.} {\bf 2024}, {\em 70},~185--202.
\newblock {\url{https://doi.org/10.17707/AgricultForest.70.3.13}}.

\bibitem[Imtiaz et~al.(2024)Imtiaz, Farooque, Randhawa, Wang, Esau, Acharya,
  and {Hashemi Garmdareh}]{IMTIAZ2024109172}
Imtiaz, F.; Farooque, A.A.; Randhawa, G.S.; Wang, X.; Esau, T.J.; Acharya, B.;
  {Hashemi Garmdareh}, S.E.
\newblock An inclusive approach to crop soil moisture estimation: Leveraging
  satellite thermal infrared bands and vegetation indices on Google Earth
  engine.
\newblock {\em Agric. Water Manag.} {\bf 2024}, {\em 306},~109172.
\newblock {\url{https://doi.org/10.1016/j.agwat.2024.109172}}.

\bibitem[Johansson and Strömquist(1978)]{JOHANSSON01011978}
Johansson, D.; Strömquist, L.
\newblock Interpretation of geomorphology and vegetation of LANDSAT satellite
  images from semi-arid central Tanzania.
\newblock {\em Nor. Geogr. Tidsskr.-Nor. J. Geogr.}
  {\bf 1978}, {\em 32},~49--54.
\newblock {\url{https://doi.org/10.1080/00291957808552025}}.

\bibitem[Cooley et~al.(2017)Cooley, Smith, Stepan, and Mascaro]{rs9121306}
Cooley, S.W.; Smith, L.C.; Stepan, L.; Mascaro, J.
\newblock Tracking Dynamic Northern Surface Water Changes with High-Frequency
  Planet CubeSat Imagery.
\newblock {\em Remote. Sens.} {\bf 2017}, {\em 9}, 1306.
\newblock {\url{https://doi.org/10.3390/rs9121306}}.

\bibitem[Reda~Sahrane and Bounab(2022)]{Sahrane13122022}
Reda~Sahrane, Y.E.K.; Bounab, A.
\newblock Investigating the effects of landscape characteristics on landslide
  susceptibility and frequency-area distributions: the case of Taounate
  province, Northern Morocco.
\newblock {\em Geocarto Int.} {\bf 2022}, {\em 37},~17686--17712.
\newblock {\url{https://doi.org/10.1080/10106049.2022.2134462}}.

\bibitem[Marco et~al.(2021)Marco, Vidal-Matutano, Morales, andAlessandro Potì,
  Kehl, Linstädter, Weniger, and Mikdad]{Marco04032021}
Marco, Y.C.; Vidal-Matutano, P.; Morales, J.; andAlessandro Potì, P.H.V.;
  Kehl, M.; Linstädter, J.; Weniger, G.C.; Mikdad, A.
\newblock Late Glacial Landscape Dynamics Based on Macrobotanical Data:
  Evidence From Ifri El Baroud (NE Morocco).
\newblock {\em Environ. Archaeol.} {\bf 2021}, {\em 26},~131--145.
\newblock {\url{https://doi.org/10.1080/14614103.2018.1538088}}.

\bibitem[Aouraghe et~al.(2010)Aouraghe, Agusti, Ouchaou, Bailon, Lopez-Garcia,
  Haddoumi, Hammouti, Oujaa, and Bougariane]{Aouraghe01032010}
Aouraghe, H.; Agusti, J.; Ouchaou, B.; Bailon, S.; Lopez-Garcia, J.M.;
  Haddoumi, H.; Hammouti, K.E.; Oujaa, A.; Bougariane, B.
\newblock The Holocene vertebrate fauna from Guenfouda site, Eastern Morocco.
\newblock {\em Hist. Biol.} {\bf 2010}, {\em 22},~320--326.
\newblock {\url{https://doi.org/10.1080/08912961003701193}}.

\bibitem[Benvenuti et~al.(2024)Benvenuti, Nesi, Papini, Risaliti, Sani, and
  Moratti]{Benvenuti31122024}
Benvenuti, M.; Nesi, J.; Papini, M.; Risaliti, G.; Sani, F.; Moratti, G.
\newblock Geology of the Toundoute Region (South Morocco): A window on the
  Early Jurassic-Cretaceous tectono-sedimentary evolution of the Central High
  Atlas.
\newblock {\em J. Maps} {\bf 2024}, {\em 20},~2419456.
\newblock {\url{https://doi.org/10.1080/17445647.2024.2419456}}.

\bibitem[Ouallali et~al.(2020)Ouallali, Aassoumi, Moukhchane, andMhammad
  Houssni, Spalevic, and Keesstra]{Ouallali03072020}
Ouallali, A.; Aassoumi, H.; Moukhchane, M.; andMhammad Houssni, A.M.; Spalevic,
  V.; Keesstra, S.
\newblock Sediment mobilization study on Cretaceous, Tertiary and Quaternary
  lithological formations of an external Rif catchment, Morocco.
\newblock {\em Hydrol. Sci. J.} {\bf 2020}, {\em 65},~1568--1582.
\newblock {\url{https://doi.org/10.1080/02626667.2020.1755435}}.

\bibitem[Bouziane~Khalloufi and Lelièvre(2010)]{Khalloufi01032010}
Bouziane~Khalloufi, D.O.; Lelièvre, H.
\newblock New paleontological and geological data about Jbel Tselfat (Late
  Cretaceous of Morocco).
\newblock {\em Hist. Biol.} {\bf 2010}, {\em 22},~57--70.
\newblock {\url{https://doi.org/10.1080/08912961003668756}}.

\bibitem[{Ali EL-Omairi} et~al.(2024){Ali EL-Omairi}, {El Garouani}, and
  Shebl]{ALIELOMAIRI2024101321}
{Ali EL-Omairi}, M.; {El Garouani}, A.; Shebl, A.
\newblock Investigation of lineament extraction: Analysis and comparison of
  digital elevation models in the Ait Semgane region, Morocco.
\newblock {\em Remote. Sens. Appl. Soc. Environ.} {\bf
  2024}, {\em 36},~101321.
\newblock {\url{https://doi.org/10.1016/j.rsase.2024.101321}}.

\bibitem[Cai et~al.(2023)Cai, Shi, Xu, and Liu]{CAI2023103226}
Cai, Y.; Shi, Q.; Xu, X.; Liu, X.
\newblock A novel approach towards continuous monitoring of forest change
  dynamics in fragmented landscapes using time series Landsat imagery.
\newblock {\em Int. J. Appl. Earth Obs. Geoinf.} {\bf 2023}, {\em 118},~103226.
\newblock {\url{https://doi.org/10.1016/j.jag.2023.103226}}.

\bibitem[Lemenkova(2024)]{Lemenkova202405}
Lemenkova, P.
\newblock {Exploitation d'images satellitaires Landsat de la région du Cap
  (Afrique du Sud) pour le calcul et la cartographie d'indices de végétation
  à l'aide du logiciel GRASS GIS}.
\newblock {\em Physio-Géo} {\bf 2024}, {\em 20},~113--129.
\newblock {\url{https://doi.org/10.4000/11pyj}}.

\bibitem[Xiong et~al.(2022)Xiong, Kang, Xu, Huang, Dai, Wang, Lu, and
  Kou]{XIONG2022109132}
Xiong, Y.; Kang, Q.; Xu, W.; Huang, S.; Dai, F.; Wang, L.; Lu, N.; Kou, W.
\newblock The dynamics of tea plantation encroachment into forests and effect
  on forest landscape pattern during 1991–2021 through time series Landsat
  images.
\newblock {\em Ecol. Indic.} {\bf 2022}, {\em 141},~109132.
\newblock {\url{https://doi.org/10.1016/j.ecolind.2022.109132}}.

\bibitem[Noy et~al.(2023)Noy, Silver, Pesek, Yizhaq, Marais, and
  Karnieli]{NOY2023103377}
Noy, K.; Silver, M.; Pesek, O.; Yizhaq, H.; Marais, E.; Karnieli, A.
\newblock Spatial and spectral analysis of fairy circles in Namibia on a
  landscape scale using satellite image processing and machine learning
  analysis.
\newblock {\em Int. J. Appl. Earth Obs. Geoinf.} {\bf 2023}, {\em 121},~103377.
\newblock {\url{https://doi.org/10.1016/j.jag.2023.103377}}.

\bibitem[Lemenkova(2024)]{Lemenkova202403}
Lemenkova, P.
\newblock {Artificial Neural Networks for Mapping Coastal Lagoon of Chilika
  Lake, India, Using Earth Observation Data}.
\newblock {\em J. Mar. Sci. Eng.} {\bf 2024}, {\em
  12},~709.
\newblock {\url{https://doi.org/10.3390/jmse12050709}}.

\bibitem[Wei~Su and Zhang(2020)]{Su02062020}
Wei~Su, D.S.; Zhang, X.
\newblock Satellite image analysis using crowdsourcing data for collaborative
  mapping: Current and opportunities.
\newblock {\em Int. J. Digit. Earth} {\bf 2020}, {\em
  13},~645--660.
\newblock {\url{https://doi.org/10.1080/17538947.2018.1556352}}.

\bibitem[Potter(2014)]{Potter18102014}
Potter, C.
\newblock Ten years of forest cover change in the Sierra Nevada detected using
  Landsat satellite image analysis.
\newblock {\em Int. J. Remote. Sens.} {\bf 2014}, {\em
  35},~7136--7153.
\newblock {\url{https://doi.org/10.1080/01431161.2014.968687}}.

\bibitem[Lemenkova(2024)]{Lemenkova202402}
Lemenkova, P.
\newblock {Deep Learning Methods of Satellite Image Processing for Monitoring
  of Flood Dynamics in the Ganges Delta, Bangladesh}.
\newblock {\em Water} {\bf 2024}, {\em 16},~1141.
\newblock {\url{https://doi.org/10.3390/w16081141}}.

\bibitem[Wang et~al.(2024)Wang, Han, Liu, Bai, Ao, Hu, Li, Li, Guo, and
  Wei]{rs16244812}
Wang, F.; Han, L.; Liu, L.; Bai, C.; Ao, J.; Hu, H.; Li, R.; Li, X.; Guo, X.;
  Wei, Y.
\newblock Advancements and Perspective in the Quantitative Assessment of Soil
  Salinity Utilizing Remote Sensing and Machine Learning Algorithms: A Review.
\newblock {\em Remote. Sens.} {\bf 2024}, {\em 16}, 4812.
\newblock {\url{https://doi.org/10.3390/rs16244812}}.

\bibitem[Cui et~al.(2025)Cui, Zhang, Du, and Li]{w17010050}
Cui, J.; Zhang, X.; Du, C.; Li, G.
\newblock Remote Sensing Identification of Harmful Algae in Ulansuhai Lake with
  Machine Learning.
\newblock {\em Water} {\bf 2025}, {\em 17}, 50.
\newblock {\url{https://doi.org/10.3390/w17010050}}.

\bibitem[Vasconcelos et~al.(2025)Vasconcelos, de~Santana, Costa, Duverger,
  Ferreira-Ferreira, Oliveira, Barbosa, Cordeiro, and
  Franca~Rocha]{fire8010008}
Vasconcelos, R.N.; de~Santana, M.M.M.; Costa, D.P.; Duverger, S.G.;
  Ferreira-Ferreira, J.; Oliveira, M.; Barbosa, L.d.S.; Cordeiro, C.L.;
  Franca~Rocha, W.J.S.
\newblock Machine Learning Model Reveals Land Use and Climate's Role in
  Caatinga Wildfires: Present and Future Scenarios.
\newblock {\em Fire} {\bf 2025}, {\em 8}, 8.
\newblock {\url{https://doi.org/10.3390/fire8010008}}.

\bibitem[Farah and Ahmed(2010)]{Farah05012010}
Farah, I.R.; Ahmed, M.B.
\newblock Towards an intelligent multi-sensor satellite image analysis based on
  blind source separation using multi-source image fusion.
\newblock {\em Int. J. Remote. Sens.} {\bf 2010}, {\em
  31},~13--38.
\newblock {\url{https://doi.org/10.1080/01431160902882504}}.

\bibitem[Lemenkova(2023)]{Lemenkova202322}
Lemenkova, P.
\newblock {Monitoring Seasonal Fluctuations in Saline Lakes of Tunisia Using
  Earth Observation Data Processed by GRASS GIS}.
\newblock {\em Land} {\bf 2023}, {\em 12},~1995.
\newblock {\url{https://doi.org/10.3390/land12111995}}.

\bibitem[Layati et~al.(2022)Layati, Ouigmane, de~Carvalho~Alves, Bagyaraj,
  Qadem, and Ghachi]{Layati03042022}
Layati, E.; Ouigmane, A.; de~Carvalho~Alves, M.; Bagyaraj, M.; Qadem, A.;
  Ghachi, M.E.
\newblock Contribution of GIS and Remote Sensing to Study and Analyze the
  Physiographic Characteristics of the Oued El-Abid Watershed (Central High
  Atlas, Morocco).
\newblock {\em Pap. Appl. Geogr.} {\bf 2022}, {\em 8},~146--162.
\newblock {\url{https://doi.org/10.1080/23754931.2021.1973545}}.

\bibitem[Layati and Ghachi(2024)]{Layati28082024}
Layati, E.; Ghachi, M.E.
\newblock Oued Lakhdar watershed (Morocco), monitoring land use/cover changes:
  remote sensing and GIS approach.
\newblock {\em Geol. Ecol. Landscapes} {\bf 2024}, {\em 0},~1--13.
\newblock {\url{https://doi.org/10.1080/24749508.2024.2395204}}.

\bibitem[{da Paz Gomes Brand{\~a}o Ferraz} and Vicens(2024)]{DAPAZGOMES}
{da Paz Gomes Brand{\~a}o Ferraz}, D.; Vicens, R.S.
\newblock Comparison between machine learning classification and
  trajectory-based change detection for identifying eucalyptus areas in Landsat
  time series.
\newblock {\em Remote. Sens. Appl. Soc. Environ.} {\bf
  2024}, p. 101444.
\newblock {\url{https://doi.org/10.1016/j.rsase.2024.101444}}.

\bibitem[Lemenkova(2020)]{Lemenkova202026}
Lemenkova, P.
\newblock {GRASS GIS for classification of Landsat TM images by maximum
  likelihood discriminant analysis: Tokyo area, Japan}.
\newblock {\em Geod. Glas.} {\bf 2020}, {\em 51},~5--25.
\newblock {\url{https://doi.org/10.6084/m9.figshare.13507761}}.

\bibitem[Laroche-Pinel et~al.(2024)Laroche-Pinel, Cianciola, Singh, Vivaldi,
  and Brillante]{LAROCHEPINEL2024109163}
Laroche-Pinel, E.; Cianciola, V.; Singh, K.; Vivaldi, G.A.; Brillante, L.
\newblock Assessing the spatial-temporal performance of machine learning in
  predicting grapevine water status from Landsat 8 imagery via block-out and
  date-out cross-validation.
\newblock {\em Agric. Water Manag.} {\bf 2024}, {\em 306},~109163.
\newblock {\url{https://doi.org/10.1016/j.agwat.2024.109163}}.

\bibitem[Nilkamal~More and Banerjee(2020)]{More02072020}
Nilkamal~More, V.B.N.; Banerjee, B.
\newblock Machine learning on high performance computing for urban greenspace
  change detection: Satellite image data fusion approach.
\newblock {\em Int. J. Image Data Fusion} {\bf 2020}, {\em
  11},~218--232.
\newblock {\url{https://doi.org/10.1080/19479832.2020.1749142}}.

\bibitem[García-Pedrero et~al.(2017)García-Pedrero, Gonzalo-Martín, and
  Lillo-Saavedra]{GarciaPedrero03042017}
García-Pedrero, A.; Gonzalo-Martín, C.; Lillo-Saavedra, M.
\newblock A machine learning approach for agricultural parcel delineation
  through agglomerative segmentation.
\newblock {\em Int. J. Remote. Sens.} {\bf 2017}, {\em
  38},~1809--1819.
\newblock {\url{https://doi.org/10.1080/01431161.2016.1278312}}.

\bibitem[Nasiri et~al.(2023)Nasiri, Hawryło, Janiec, and
  Socha]{NASIRI2023103555}
Nasiri, V.; Hawryło, P.; Janiec, P.; Socha, J.
\newblock Comparing Object-Based and Pixel-Based Machine Learning Models for
  Tree-Cutting Detection with PlanetScope Satellite Images: Exploring Model
  Generalization.
\newblock {\em Int. J. Appl. Earth Obs. Geoinf.} {\bf 2023}, {\em 125},~103555.
\newblock {\url{https://doi.org/10.1016/j.jag.2023.103555}}.

\bibitem[Ebrahimy and Zhang(2023)]{EBRAHIMY2023103390}
Ebrahimy, H.; Zhang, Z.
\newblock Per-pixel accuracy as a weighting criterion for combining ensemble of
  extreme learning machine classifiers for satellite image classification.
\newblock {\em Int. J. Appl. Earth Obs. Geoinf.} {\bf 2023}, {\em 122},~103390.
\newblock {\url{https://doi.org/10.1016/j.jag.2023.103390}}.

\bibitem[Fu et~al.(2024)Fu, Li, Fu, Wang, and Wang]{FU2024104421}
Fu, H.; Li, H.; Fu, A.; Wang, X.; Wang, Q.
\newblock Transportation emissions monitoring and policy research: Integrating
  machine learning and satellite imaging.
\newblock {\em Transp. Res. Part Transp. Environ.} {\bf
  2024}, {\em 136},~104421.
\newblock {\url{https://doi.org/10.1016/j.trd.2024.104421}}.

\bibitem[Padhy et~al.(2024)Padhy, Dash, and Mishra]{PADHY2024100278}
Padhy, R.; Dash, S.K.; Mishra, J.
\newblock Joint feature selection and classification of low-resolution
  satellite images using the SAT-6 dataset.
\newblock {\em High-Confid. Comput.} {\bf 2024}, p. 100278.
\newblock {\url{https://doi.org/10.1016/j.hcc.2024.100278}}.

\bibitem[Lemenkova(2024)]{Lemenkova202401}
Lemenkova, P.
\newblock {Random Forest Classifier Algorithm of Geographic Resources Analysis
  Support System Geographic Information System for Satellite Image Processing:
  Case Study of Bight of Sofala, Mozambique}.
\newblock {\em Coasts} {\bf 2024}, {\em 4},~127--149.
\newblock {\url{https://doi.org/10.3390/coasts4010008}}.

\bibitem[Habib and Okayli(2023)]{HABIB20231}
Habib, M.; Okayli, M.
\newblock An Overview of Modern Cartographic Trends Aligned with the ICA’s
  Perspective.
\newblock {\em Rev. Int. Geomat.} {\bf 2023}, {\em 32},~1--16.
\newblock {\url{https://doi.org/10.32604/RIG.2023.043399}}.

\bibitem[Lemenkova(2024)]{Lemenkova202406}
Lemenkova, P.
\newblock {Landscape Fragmentation and Deforestation in Sierra Leone, West
  Africa, Analysed Using Satellite Images}.
\newblock {\em Transylv. Rev. Syst. Ecol. Res.}
  {\bf 2024}, {\em 26},~13--26.
\newblock {\url{https://doi.org/10.2478/trser-2024-0002}}.

\bibitem[Khaldi et~al.(2024)Khaldi, Tabik, Puertas-Ruiz, {de Giles}, Correa,
  Zamora, and Segura]{KHALDI2024104191}
Khaldi, R.; Tabik, S.; Puertas-Ruiz, S.; {de Giles}, J.P.; Correa, J.A.H.;
  Zamora, R.; Segura, D.A.
\newblock Individual mapping of large polymorphic shrubs in high mountains
  using satellite images and deep learning.
\newblock {\em Int. J. Appl. Earth Obs. Geoinf.} {\bf 2024}, {\em 134},~104191.
\newblock {\url{https://doi.org/10.1016/j.jag.2024.104191}}.

\bibitem[Lemenkova and Debeir(2022)]{Lemenkova202227}
Lemenkova, P.; Debeir, O.
\newblock {Satellite Image Processing by Python and R Using Landsat 9 OLI/TIRS
  and SRTM DEM Data on Côte d’Ivoire, West Africa}.
\newblock {\em J. Imaging} {\bf 2022}, {\em 8},~317.
\newblock {\url{https://doi.org/10.3390/jimaging8120317}}.

\bibitem[Pereira et~al.(2023)Pereira, Pereira, Gil, and
  Mantas]{PEREIRA2023106653}
Pereira, J.; Pereira, A.; Gil, A.; Mantas, V.M.
\newblock Lithology mapping with satellite images, fieldwork-based spectral
  data, and machine learning algorithms: The case study of Beiras Group
  (Central Portugal).
\newblock {\em CATENA} {\bf 2023}, {\em 220},~106653.
\newblock {\url{https://doi.org/10.1016/j.catena.2022.106653}}.

\bibitem[Eskandari and Bordbar(2024)]{ESKANDARI2024}
Eskandari, S.; Bordbar, S.K.
\newblock Provision of land use and forest density maps in semi-arid areas of
  Iran using Sentinel-2 satellite images and vegetation indices.
\newblock {\em Adv. Space Res.} {\bf \hl{2024}}.
\newblock {\url{https://doi.org/10.1016/j.asr.2024.10.060}}.
%MDPI: Please add the volume and page number.


\bibitem[Aboutofail and Slimani(2024)]{ABOUTOFAIL2024112522}
Aboutofail, S.; Slimani, H.
\newblock The Paleocene--Eocene Thermal Maximum (PETM) interval in the
  southwestern Mediterranean Tethys at Morocco: New data from a high-resolution
  study of dinoflagellate cysts and palynofacies in the Rif Chain.
\newblock {\em Palaeogeogr. Palaeoclimatol. Palaeoecol.} {\bf 2024},
  {\em 655},~112522.

\bibitem[McGregor et~al.(2009)McGregor, Dupont, Stuut, and
  Kuhlmann]{MCGREGOR20091434}
McGregor, H.V.; Dupont, L.; Stuut, J.B.W.; Kuhlmann, H.
\newblock {Vegetation change, goats, and religion: a 2000-year history of land
  use in southern Morocco}.
\newblock {\em Quat. Sci. Rev.} {\bf 2009}, {\em 28},~1434--1448.
\newblock {\url{https://doi.org/10.1016/j.quascirev.2009.02.012}}.

\bibitem[Criniti et~al.(2024)Criniti, Martín-Martín, Hlila, Maaté, and
  Maaté]{CRINITI2024106861}
Criniti, S.; Martín-Martín, M.; Hlila, R.; Maaté, A.; Maaté, S.
\newblock Detrital signatures of the Ghomaride Culm cycle (Rif Cordillera, N
  Morocco): New constraints for the northern Gondwana plate tectonics.
\newblock {\em Mar. Pet. Geol.} {\bf 2024}, {\em 165},~106861.
\newblock {\url{https://doi.org/10.1016/j.marpetgeo.2024.106861}}.

\bibitem[Parish and Funnell(1999)]{PARISH199945}
Parish, R.; Funnell, D.
\newblock Climate change in mountain regions: some possible consequences in the
  Moroccan High Atlas.
\newblock {\em Glob. Environ. Chang.} {\bf 1999}, {\em 9},~45--58.
\newblock {\url{https://doi.org/10.1016/S0959-3780(98)00021-1}}.

\bibitem[{El Talibi} et~al.(2014){El Talibi}, Zaghloul, Perri, Aboumaria,
  Rossi, and {El Moussaoui}]{ELTALIBI2014554}
{El Talibi}, H.; Zaghloul, M.N.; Perri, F.; Aboumaria, K.; Rossi, A.; {El
  Moussaoui}, S.
\newblock Sedimentary evolution of the siliciclastic Aptian--Albian Massylian
  flysch of the Chouamat Nappe (central Rif, Morocco).
\newblock {\em J. Afr. Earth Sci.} {\bf 2014}, {\em
  100},~554--568.
\newblock {\url{https://doi.org/10.1016/j.jafrearsci.2014.08.004}}.

\bibitem[Bonardi et~al.(2003)Bonardi, {de Capoa}, {Di Staso}, Est{\'e}vez,
  Mart{\'\i}n-Mart{\'\i}n, Mart{\'\i}n-Rojas, Perrone, and
  Tent-Mancl{\'u}s]{BONARDI2003149}
Bonardi, G.; {de Capoa}, P.; {Di Staso}, A.; Est{\'e}vez, A.;
  Mart{\'\i}n-Mart{\'\i}n, M.; Mart{\'\i}n-Rojas, I.; Perrone, V.;
  Tent-Mancl{\'u}s, J.E.
\newblock Oligocene-to-Early Miocene depositional and structural evolution of
  the Calabria--Peloritani Arc southern terrane (Italy) and geodynamic
  correlations with the Spain Betics and Morocco Rif.
\newblock {\em Geodin. Acta} {\bf 2003}, {\em 16},~149--169.
\newblock {\url{https://doi.org/10.1016/j.geoact.2003.06.001}}.

\bibitem[Maat{\'e} et~al.(2000)Maat{\'e}, Mart{\'\i}n-Algarra,
  Mart{\'\i}n-Mart{\'\i}n, and Serra-Kiel]{MAATE2000409}
Maat{\'e}, A.; Mart{\'\i}n-Algarra, A.; Mart{\'\i}n-Mart{\'\i}n, M.;
  Serra-Kiel, J.
\newblock Nouvelles donn{\'e}essur le Pal{\'e}oc{\`e}ne-{\'E}oc{\`e}ne des
  zones internes b{\'e}tico-rifaines.
\newblock {\em Geobios} {\bf 2000}, {\em 33},~409--418.
\newblock {\url{https://doi.org/10.1016/S0016-6995(00)80074-1}}.

\bibitem[Slimani et~al.(2019)Slimani, Mahboub, Toufiq, Jbari, Chakir, and
  Tahiri]{SLIMANI2019101785}
Slimani, H.; Mahboub, I.; Toufiq, A.; Jbari, H.; Chakir, S.; Tahiri, A.
\newblock Bartonian to Priabonian dinoflagellate cyst biostratigraphy and
  paleoenvironments of the M'karcha section in the Southern Tethys margin (Rif
  Chain, Northern Morocco).
\newblock {\em Mar. Micropaleontol.} {\bf 2019}, {\em 153},~101785.
\newblock {\url{https://doi.org/10.1016/j.marmicro.2019.101785}}.

\bibitem[Mart{\'\i}n-Mart{\'\i}n et~al.(2023)Mart{\'\i}n-Mart{\'\i}n, Guerrera,
  Ca{\~n}averas, Alcal{\'a}, Serrano, Maat{\'e}, Hlila, Maat{\'e}, Tramontana,
  S{\'a}nchez-Navas, and {Le Breton}]{MARTINMARTIN2023106367}
Mart{\'\i}n-Mart{\'\i}n, M.; Guerrera, F.; Ca{\~n}averas, J.C.; Alcal{\'a},
  F.J.; Serrano, F.; Maat{\'e}, A.; Hlila, R.; Maat{\'e}, S.; Tramontana, M.;
  S{\'a}nchez-Navas, A.;  et~al.
\newblock Paleogene evolution of the External Rif Zone (Morocco) and comparison
  with other western Tethyan margins.
\newblock {\em Sediment. Geol.} {\bf 2023}, {\em 448},~106367.
\newblock {\url{https://doi.org/10.1016/j.sedgeo.2023.106367}}.

\bibitem[{Ait Brahim} et~al.(2002){Ait Brahim}, Chotin, Hinaj, Abdelouafi, {El
  Adraoui}, Nakcha, Dhont, Charroud, {Sossey Alaoui}, Amrhar, Bouaza, Tabyaoui,
  and Chaouni]{AITBRAHIM2002187}
{Ait Brahim}, L.; Chotin, P.; Hinaj, S.; Abdelouafi, A.; {El Adraoui}, A.;
  Nakcha, C.; Dhont, D.; Charroud, M.; {Sossey Alaoui}, F.; Amrhar, M.;  et~al.
\newblock Paleostress evolution in the Moroccan African margin from Triassic to
  Present.
\newblock {\em Tectonophysics} {\bf 2002}, {\em 357},~187--205.
\newblock {\url{https://doi.org/10.1016/S0040-1951(02)00368-2}}.

\bibitem[Leblanc and Olivier(1984)]{LEBLANC1984345}
Leblanc, D.; Olivier, P.
\newblock Role of strike-slip faults in the Betic-Rifian orogeny.
\newblock {\em Tectonophysics} {\bf 1984}, {\em 101},~345--355.
\newblock {\url{https://doi.org/0.1016/0040-1951(84)90120-3}}.

\bibitem[Ara{\~n}a and Vegas(1974)]{ARANA1974197}
Ara{\~n}a, V.; Vegas, R.
\newblock Plate tectonics and volcanism in the gibraltar arc.
\newblock {\em Tectonophysics} {\bf 1974}, {\em 24},~197--212.
\newblock {\url{https://doi.org/10.1016/0040-1951(74)90008-0}}.

\bibitem[Perri et~al.(2022)Perri, Mart{\'\i}n-Mart{\'\i}n, Maat{\'e}, Hlila,
  Maat{\'e}, Criniti, Capobianco, and Critelli]{PERRI2022105811}
Perri, F.; Mart{\'\i}n-Mart{\'\i}n, M.; Maat{\'e}, A.; Hlila, R.; Maat{\'e},
  S.; Criniti, S.; Capobianco, W.; Critelli, S.
\newblock Provenance and paleogeographic implications for the Cenozoic
  sedimentary cover of the Ghomaride Complex (Internal Rif Belt), Morocco.
\newblock {\em Mar. Pet. Geol.} {\bf 2022}, {\em 143},~105811.
\newblock {\url{https://doi.org/10.1016/j.marpetgeo.2022.105811}}.

\bibitem[Talmat et~al.(2025)Talmat, Mart{\'\i}n-Mart{\'\i}n, Benyoucef,
  Ferr{\'e}, and Belhai]{TALMAT2025105515}
Talmat, S.; Mart{\'\i}n-Mart{\'\i}n, M.; Benyoucef, M.; Ferr{\'e}, B.; Belhai,
  D.
\newblock The Koudiat El Madene unit (Kabylian ``Dorsal'', Algeria) and its
  correlation with similar units belonging to the Rif-Betic Arc (Morocco and
  Spain).
\newblock {\em J. Afr. Earth Sci.} {\bf 2025}, {\em
  223},~105515.
\newblock {\url{https://doi.org/10.1016/j.jafrearsci.2024.105515}}.

\bibitem[Chakkour et~al.(2024)Chakkour, Ennouni, Sahli, Kadaoui, Houssni,
  Kassout, and Ater]{CHAKKOUR2024}
Chakkour, S.; Ennouni, H.; Sahli, A.; Kadaoui, K.; Houssni, M.; Kassout, J.;
  Ater, M.
\newblock Vegetation composition and ecological characteristics of the
  fragmented Alnus glutinosa woodlands of Morocco.
\newblock {\em Ecol. Front.}{  \bf  2024}, \emph{45}, 54--67.
\newblock {\url{https://doi.org/10.1016/j.ecofro.2024.08.007}}.




\bibitem[Hamouch et~al.(2024)Hamouch, Chaaouan, and eddine Bouiss]{HAMOUCH2024}
Hamouch, C.; Chaaouan, J.; eddine Bouiss, C.
\newblock Spatial-temporal assessment of soil erosion using the RUSLE model in
  the upstream Inaou{\`e}ne watershed, Northern Morocco.
\newblock {\em Nat. Hazards Res.} {\bf \hl{2024}}.
\newblock {\url{https://doi.org/10.1016/j.nhres.2024.08.002}}.
%MDPI: Please add the volume and page number.



\bibitem[Ed-Dakiri et~al.(2024)Ed-Dakiri, Etebaai, Moussaoui, Tawfik,
  Lamgharbaj, Talibi, Dekkaki, and Taher]{EDDAKIRI2024e02401}
Ed-Dakiri, S.; Etebaai, I.; Moussaoui, S.E.; Tawfik, A.; Lamgharbaj, M.;
  Talibi, H.E.; Dekkaki, H.C.; Taher, M.
\newblock Assessing soil erosion risk through geospatial analysis and magnetic
  susceptibility: A study in the Oued Ghiss dam watershed, Central Rif,
  Morocco.
\newblock {\em Sci. Afr.} {\bf 2024}, {\em 26},~e02401.
\newblock {\url{https://doi.org/10.1016/j.sciaf.2024.e02401}}.

\bibitem[Boubekraoui et~al.(2023)Boubekraoui, Maouni, Ghallab, Draoui, and
  Maouni]{BOUBEKRAOUI2023100388}
Boubekraoui, H.; Maouni, Y.; Ghallab, A.; Draoui, M.; Maouni, A.
\newblock Spatio-temporal analysis and identification of deforestation hotspots
  in the Moroccan western Rif.
\newblock {\em Trees For. People} {\bf 2023}, {\em 12},~100388.
\newblock {\url{https://doi.org/10.1016/j.tfp.2023.100388}}.

\bibitem[Aqil et~al.(2015)Aqil, Khattach, Gout, Mandour, and
  Kaufmann]{Aqil03042015}
Aqil, H.; Khattach, D.; Gout, R.E.; Mandour, A.E.; Kaufmann, O.
\newblock Contribution de la gravimétrie à l’étude de la structure
  profonde du bassin de Bou-Houria (Maroc Nord-oriental) : implications
  hydrogéologiques.
\newblock {\em Hydrol. Sci. J.} {\bf 2015}, {\em 60},~736--745.
\newblock {\url{https://doi.org/10.1080/02626667.2014.943229}}.

\bibitem[Zouhri et~al.(2004)Zouhri, Gorini, Mania, Deffontaines, and
  Zerouali]{Zouhri01062004}
Zouhri, L.; Gorini, C.; Mania, J.; Deffontaines, B.; Zerouali, A.E.H.
\newblock Spatial distribution of resistivity in the hydrogeological systems,
  and identification of the catchment area in the Rharb basin, Morocco /
  Répartition spatiale de la résistivité dans les systèmes
  hydrogéologiques et détection des zones de captages dans le bassin du
  Rharb, Maroc.
\newblock {\em Hydrol. Sci. J.} {\bf 2004}, {\em 49},~--398.
\newblock {\url{https://doi.org/10.1623/hysj.49.3.387.54350}}.

\bibitem[Al~Karkouri(2017)]{al2017milieux}
Al~Karkouri, J.
\newblock Les milieux montagneux marocains {\`a} l’{\'e}preuve du changement
  climatique (cas de la montagne rifaine).
\newblock {\em Hesp{\'e}ris-Tamuda} {\bf 2017}, {\em 52},~237--267.

\bibitem[Duarte et~al.(2024)Duarte, Magalhães, Hernández-Molina, Roque, and
  Menapace]{DUARTE2024104330}
Duarte, D.; Magalhães, V.H.; Hernández-Molina, F.J.; Roque, C.; Menapace, W.
\newblock Interaction between active tectonics, bottom-current processes and
  coral mounds: A unique example in the NW Moroccan Margin, southern Gulf of
  Cadiz.
\newblock {\em Deep. Sea Res. Part I Oceanogr. Res. Pap.} {\bf
  2024}, {\em 209},~104330.
\newblock {\url{https://doi.org/10.1016/j.dsr.2024.104330}}.

\bibitem[Chabli et~al.(2014)Chabli, Chalouan, Akil, Galindo-Zaldívar, Ruano,
  {Sanz de Galdeano}, López-Garrido, Marín-Lechado, and
  Pedrera]{CHABLI2014123}
Chabli, A.; Chalouan, A.; Akil, M.; Galindo-Zaldívar, J.; Ruano, P.; {Sanz de
  Galdeano}, C.; López-Garrido, A.C.; Marín-Lechado, C.; Pedrera, A.
\newblock Plio-Quaternary paleostresses in the Atlantic passive margin of the
  Moroccan Meseta: Influence of the Central Rif escape tectonics related to
  Eurasian-African plate convergence.
\newblock {\em J. Geodyn.} {\bf 2014}, {\em 77},~123--134.
\newblock SI : Geodynamic evolution of the Alboran domain,
  {\url{https://doi.org/10.1016/j.jog.2013.12.005}}.

\bibitem[Lghamour et~al.(2024)Lghamour, Karrat, Picotti, Hajdas, Haghipour,
  Guidobaldi, {Wyss Heeb}, and Gourari]{LGHAMOUR2024108941}
Lghamour, M.; Karrat, L.; Picotti, V.; Hajdas, I.; Haghipour, N.; Guidobaldi,
  G.; {Wyss Heeb}, K.; Gourari, L.
\newblock Late Pleistocene to Holocene alluvial deposits of the Inaouène
  Valley and their paleoenvironmental significance (north Morocco).
\newblock {\em Quat. Sci. Rev.} {\bf 2024}, {\em 344},~108941.
\newblock {\url{https://doi.org/10.1016/j.quascirev.2024.108941}}.

\bibitem[Kairouani et~al.(2024)Kairouani, Abbassi, Zaghloul, {El Mourabet},
  Micheletti, Fornelli, Mongelli, and Critelli]{KAIROUANI2024106762}
Kairouani, H.; Abbassi, A.; Zaghloul, M.N.; {El Mourabet}, M.; Micheletti, F.;
  Fornelli, A.; Mongelli, G.; Critelli, S.
\newblock The Jurassic climate change in the northwest Gondwana (External Rif,
  Morocco): Evidence from geochemistry and implication for paleoclimate
  evolution.
\newblock {\em Mar. Pet. Geol.} {\bf 2024}, {\em 163},~106762.
\newblock {\url{https://doi.org/10.1016/j.marpetgeo.2024.106762}}.

\bibitem[Malusà et~al.(2024)Malusà, Ellero, and Ottria]{MALUSA2024230533}
Malusà, M.G.; Ellero, A.; Ottria, G.
\newblock Tectonics of the Mw 6.8 Al Haouz earthquake (Morocco) reveals minor
  role of asthenospheric upwelling.
\newblock {\em Tectonophysics} {\bf 2024}, {\em 891},~230533.
\newblock {\url{https://doi.org/10.1016/j.tecto.2024.230533}}.

\bibitem[Sadkaoui et~al.(2024)Sadkaoui, Brahim, Kader, Agharroud, Mihraje,
  Aluni, Aassoumi, Sassioui, Spalevic, and Sestras]{SADKAOUI2024105219}
Sadkaoui, D.; Brahim, B.; Kader, S.; Agharroud, K.; Mihraje, A.I.; Aluni, K.;
  Aassoumi, H.; Sassioui, S.; Spalevic, V.; Sestras, P.
\newblock Evaluation of tectonic activity using morphometric indices: Study of
  the case of Taïliloute ridge (middle-Atlas region, Morocco).
\newblock {\em J. Afr. Earth Sci.} {\bf 2024}, {\em
  213},~105219.
\newblock {\url{https://doi.org/10.1016/j.jafrearsci.2024.105219}}.

\bibitem[Costa et~al.(2024)Costa, da~Conceição~Freitas, Teixeira, Costa,
  Santos, Cachão, Fatela, Pereira, Tereso, Lopes, Bao, and
  El-Boudjay]{COSTA2024109578}
Costa, A.M.; da~Conceição~Freitas, M.; Teixeira, A.; Costa, P.J.; Santos,
  R.N.; Cachão, M.; Fatela, F.; Pereira, R.; Tereso, J.P.; Lopes, V.P.;
  et~al.
\newblock Environmental changes and historical occupation at Oued Laksar (Ksar
  Seghir), Morocco.
\newblock {\em Quat. Int.} {\bf 2024}, {\em 712},~109578.
\newblock {\url{https://doi.org/10.1016/j.quaint.2024.10.010}}.

\bibitem[Boutahar et~al.(2023)Boutahar, Gonzalez, Picone, Crisafulli, Mesa,
  Redouan, {El Bakali}, Kadiri, Lamrani, and Merzouki]{BOUTAHAR2023104828}
Boutahar, A.; Gonzalez, P.C.; Picone, R.M.; Crisafulli, A.; Mesa, J.M.;
  Redouan, F.Z.; {El Bakali}, I.; Kadiri, M.; Lamrani, Z.; Merzouki, A.
\newblock Modern pollen–vegetation relationship in the Rif mountains
  (Northern Morocco).
\newblock {\em Rev. Palaeobot. Palynol.} {\bf 2023}, {\em
  310},~104828.
\newblock {\url{https://doi.org/10.1016/j.revpalbo.2022.104828}}.

\bibitem[Muller et~al.(2022)Muller, Daoud-Bouattour, Fauquette,
  Bottollier-Curtet, Rifai, Robles, Saber, {El Madihi}, Moukrim, and
  Rhazi]{MULLER202235}
Muller, S.D.; Daoud-Bouattour, A.; Fauquette, S.; Bottollier-Curtet, M.; Rifai,
  N.; Robles, M.; Saber, E.R.; {El Madihi}, M.; Moukrim, S.; Rhazi, L.
\newblock Holocene history of peatland communities of central Rif (Northern
  Morocco).
\newblock {\em Geobios} {\bf 2022}, {\em 70},~35--53.
\newblock {\url{https://doi.org/10.1016/j.geobios.2021.12.001}}.

\bibitem[Ajbilou et~al.(2006)Ajbilou, Mara{\~n}{\'o}n, and
  Arroyo]{AJBILOU2006104}
Ajbilou, R.; Mara{\~n}{\'o}n, T.; Arroyo, J.
\newblock Ecological and biogeographical analyses of Mediterranean forests of
  northern Morocco.
\newblock {\em Acta Oecologica} {\bf 2006}, {\em 29},~104--113.
\newblock {\url{https://doi.org/10.1016/j.actao.2005.08.006}}.

\bibitem[Mara{\~n}{\'o}n et~al.(1999)Mara{\~n}{\'o}n, Ajbilou, Ojeda, and
  Arroyo]{MARANON1999147}
Mara{\~n}{\'o}n, T.; Ajbilou, R.; Ojeda, F.; Arroyo, J.
\newblock Biodiversity of woody species in oak woodlands of southern Spain and
  northern Morocco.
\newblock {\em For. Ecol. Manag.} {\bf 1999}, {\em 115},~147--156.
\newblock {\url{https://doi.org/10.1016/S0378-1127(98)00395-8}}.

\bibitem[Salmon et~al.(2015)Salmon, Friedl, Frolking, Wisser, and
  Douglas]{SALMON2015321}
Salmon, J.; Friedl, M.A.; Frolking, S.; Wisser, D.; Douglas, E.M.
\newblock Global rain-fed, irrigated, and paddy croplands: A new high
  resolution map derived from remote sensing, crop inventories and climate
  data.
\newblock {\em Int. J. Appl. Earth Obs. Geoinf.} {\bf 2015}, {\em 38},~321--334.
\newblock {\url{https://doi.org/10.1016/j.jag.2015.01.014}}.

\bibitem[Zayani et~al.(2023)Zayani, Ammari, {Ben Allal}, and
  Bouhafa]{ZAYANI2023e22910}
Zayani, I.; Ammari, M.; {Ben Allal}, L.; Bouhafa, K.
\newblock Agroforestry olive orchards for soil organic carbon storage: Case of
  Saiss, Morocco.
\newblock {\em Heliyon} {\bf 2023}, {\em 9},~e22910.
\newblock {\url{https://doi.org/10.1016/j.heliyon.2023.e22910}}.

\bibitem[Lemenkova(2022)]{Lemenkova202221}
Lemenkova, P.
\newblock {Evapotranspiration, vapour pressure and climatic water deficit in
  Ethiopia mapped using GMT and TerraClimate dataset}.
\newblock {\em J. Water Land Dev.} {\bf 2022}, {\em
  54},~201--209.
\newblock {\url{https://doi.org/10.24425/jwld.2022.141573}}.

\bibitem[Becker(2005)]{BECKER20051075}
Becker, N.C.
\newblock Painting by numbers: A GMT primer for merging swath-mapping sonar
  data of different types and resolutions.
\newblock {\em Comput. Geosci.} {\bf 2005}, {\em 31},~1075--1077.
\newblock {\url{https://doi.org/10.1016/j.cageo.2005.02.016}}.

\bibitem[Lemenkova(2022)]{Lemenkova202217}
Lemenkova, P.
\newblock {Mapping Climate Parameters over the Territory of Botswana Using GMT
  and Gridded Surface Data from TerraClimate}.
\newblock {\em ISPRS Int. J.-Geo-Inf.} {\bf 2022}, {\em
  11},~473.
\newblock {\url{https://doi.org/10.3390/ijgi11090473}}.

\bibitem[Bawa et~al.(2024)Bawa, Onotu, Akomolafe, and Sa’i]{10537274}
Bawa, S.; Onotu, A.A.; Akomolafe, E.A.; Sa’i, U.I.
\newblock Understanding the Topography of the Gulf of Guinea Seabed Using GMT
  Scripting and GEBCO Gridded Data.
\newblock In Proceedings of the 2024 IEEE Mediterranean and Middle-East
  Geoscience and Remote Sensing Symposium (M2GARSS),  \hl{Oran, Algeria, 15--17 April 2024}; pp. 516--520.
\newblock {\url{https://doi.org/10.1109/M2GARSS57310.2024.10537274}}.
%MDPI: We added the location and date of the conference. Please confirm. Same as the following highlight.


\bibitem[Shi et~al.(2009)Shi, Du, Lu, Hu, and Ke]{5369602}
Shi, H.; Du, Z.; Lu, Y.; Hu, X.; Ke, X.
\newblock Amery Ice Shelf Digital Elevation Model from GLAS and GMT.
\newblock In Proceedings of the 2009 Third International Symposium on
  Intelligent Information Technology Application,  \hl{Nan Chang, China, 21--22 November 2009}; Volume~2, pp. 129--133.
\newblock {\url{https://doi.org/10.1109/IITA.2009.313}}.

\bibitem[Fang et~al.(2020)Fang, Jiang, Xu, and Wang]{FANG202031}
Fang, Z.; Jiang, G.; Xu, C.; Wang, S.
\newblock A tectonic geodesy mapping software based on QGIS.
\newblock {\em Geod. Geodyn.} {\bf 2020}, {\em 11},~31--39.
\newblock {\url{https://doi.org/10.1016/j.geog.2019.08.001}}.

\bibitem[Wang et~al.(2023)Wang, Peng, Hu, and Zhang]{f14071414}
Wang, Y.; Peng, Y.; Hu, X.; Zhang, P.
\newblock Fine-Resolution Forest Height Estimation by Integrating ICESat-2 and
  Landsat 8 OLI Data with a Spatial Downscaling Method for Aboveground Biomass
  Quantification.
\newblock {\em Forests} {\bf 2023}, {\em 14}, 1414.
\newblock {\url{https://doi.org/10.3390/f14071414}}.

\bibitem[{Abd El-Rahman Hegab} et~al.(2024){Abd El-Rahman Hegab}, {Abou El
  Magd}, and {Hamed Abd El Wahid}]{ABDELRAHMANHEGAB2024716}
{Abd El-Rahman Hegab}, M.; {Abou El Magd}, I.; {Hamed Abd El Wahid}, K.
\newblock Revealing Potential Mineralization Zones Utilizing Landsat-9, ASTER
  and Airborne Radiometric Data at Elkharaza-Dara Area, North Eastern Desert,
  Egypt.
\newblock {\em  Egypt. J. Remote. Sens. Space Sci.} {\bf
  2024}, {\em 27},~716--733.
\newblock {\url{https://doi.org/10.1016/j.ejrs.2024.10.005}}.

\bibitem[Xu et~al.(2023)Xu, Tan, Abd-Elrahman, Fan, and Wang]{rs15082017}
Xu, Y.; Tan, Y.; Abd-Elrahman, A.; Fan, T.; Wang, Q.
\newblock Incorporation of Fused Remote Sensing Imagery to Enhance Soil Organic
  Carbon Spatial Prediction in an Agricultural Area in Yellow River Basin,
  China.
\newblock {\em Remote. Sens.} {\bf 2023}, {\em 15}, 2017.
\newblock {\url{https://doi.org/10.3390/rs15082017}}.

\bibitem[Aghababaei et~al.(2021)Aghababaei, Ebrahimi, Naghipour, Asadi, and
  Verrelst]{rs13173433}
Aghababaei, M.; Ebrahimi, A.; Naghipour, A.A.; Asadi, E.; Verrelst, J.
\newblock Classification of Plant Ecological Units in Heterogeneous Semi-Steppe
  Rangelands: Performance Assessment of Four Classification Algorithms.
\newblock {\em Remote Sens.} {\bf 2021}, {\em 13}, 3433.
\newblock {\url{https://doi.org/10.3390/rs13173433}}.

\bibitem[Kusi et~al.(2021)Kusi, Khattabi, and Mhammdi]{Kusi}
Kusi, K.K.; Khattabi, A.; Mhammdi, N.
\newblock Integrated assessment of ecosystem services in response to land use
  change and management activities in Morocco.
\newblock {\em Arab. J. Geosci.} {\bf 2021}, {\em 14},~418.
\newblock {\url{https://doi.org/10.1007/s12517-021-06719-x}}.

\bibitem[Hamdani and Baali(2024)]{Hamdani}
Hamdani, N.; Baali, A., Impact of Land Use/Land Cover (LULC) Changes on the
  Watersheds of Three Lakes in the Central Middle Atlas, (Morocco).
\newblock In {\em Advanced Technology for Smart Environment and Energy};
  Springer Nature: Cham, Switzerland, 2024; \mbox{pp. 185--196.}
\newblock {\url{https://doi.org/10.1007/978-3-031-50871-4_11}}.

\bibitem[Kusi et~al.(2023)Kusi, Khattabi, and Mhammdi]{Kusi2023}
Kusi, K.K.; Khattabi, A.; Mhammdi, N.
\newblock Analyzing the impact of land use change on ecosystem service value in
  the main watersheds of Morocco.
\newblock {\em Environ. Dev. Sustain.} {\bf 2023}, {\em
  25},~2688--2715.
\newblock {\url{https://doi.org/10.1007/s10668-022-02162-4}}.

\bibitem[El~fallah et~al.(2024)El~fallah, El~kharrim, and Belghyti]{Elfallah}
El~fallah, K.; El~kharrim, K.; Belghyti, D.
\newblock Land use land cover change detection by using remote sensing in
  Meknes province, Morocco with an indicator based (DPSIR) approach.
\newblock {\em Vegetos} {\bf \hl{2024}}.
\newblock {\url{https://doi.org/10.1007/s42535-024-01110-z}}.
%MDPI: Please add the volume and page number.



\bibitem[Ido(2024)]{Idoumskine}
{\em Diachronic Analysis of Land Use/Land Cover Impacts on Carbon Storage in
  Agadir City (Morocco)}; Springer Nature: Cham,  Switzerland, 2024.  
\newblock {\url{https://doi.org/10.1007/978-3-031-43218-7_50}}.

\bibitem[Zaher et~al.(2020)Zaher, Sabir, Benjelloun, and
  Paul-Igor]{ZAHER2020109544}
Zaher, H.; Sabir, M.; Benjelloun, H.; Paul-Igor, H.
\newblock Effect of forest land use change on carbohydrates, physical soil
  quality and carbon stocks in Moroccan cedar area.
\newblock {\em J. Environ. Manag.} {\bf 2020}, {\em
  254},~109544.
\newblock {\url{https://doi.org/10.1016/j.jenvman.2019.109544}}.

\bibitem[Moumane et~al.(2022)Moumane, {Al Karkouri}, Benmansour, {El Ghazali},
  Fico, Karmaoui, and Batchi]{MOUMANE2022100745}
Moumane, A.; {Al Karkouri}, J.; Benmansour, A.; {El Ghazali}, F.E.; Fico, J.;
  Karmaoui, A.; Batchi, M.
\newblock Monitoring long-term land use, land cover change, and desertification
  in the Ternata oasis, Middle Draa Valley, Morocco.
\newblock {\em Remote. Sens. Appl. Soc. Environ.} {\bf
  2022}, {\em 26},~100745.
\newblock {\url{https://doi.org/10.1016/j.rsase.2022.100745}}.

\bibitem[Mulas et~al.(2012)Mulas, Bellavite, Lubino, Belkheiri, and
  Enne]{MULAS2012e46}
Mulas, M.; Bellavite, D.; Lubino, M.; Belkheiri, O.; Enne, G.
\newblock Participatory Approach for Integrated Development and Management of
  North African Marginal Zones: Demonstrative Plan to Fight Desertification in
  Morocco and Tunisia.
\newblock {\em Ital. J. Agron.} {\bf 2012}, {\em 7},~e46.
\newblock {\url{https://doi.org/10.4081/ija.2012.e46}}.

\bibitem[Kouba et~al.(2018)Kouba, Gartzia, {El Aich}, and Alados]{KOUBA20181}
Kouba, Y.; Gartzia, M.; {El Aich}, A.; Alados, C.L.
\newblock Deserts do not advance, they are created: Land degradation and
  desertification in semiarid environments in the Middle Atlas, Morocco.
\newblock {\em J. Arid. Environ.} {\bf 2018}, {\em 158},~1--8.
\newblock {\url{https://doi.org/10.1016/j.jaridenv.2018.07.002}}.

\bibitem[Mayaux and Fernandez(2024)]{MAYAUX2024104093}
Mayaux, P.L.; Fernandez, S.
\newblock Blinded like a state: Water scarcity and the quantification dilemma
  in Morocco.
\newblock {\em Geoforum} {\bf 2024}, {\em 155},~104093.
\newblock {\url{https://doi.org/10.1016/j.geoforum.2024.104093}}.

\bibitem[Chebli et~al.(2018)Chebli, Chentouf, Ozer, Hornick, and
  Cabaraux]{chebli2018forest}
Chebli, Y.; Chentouf, M.; Ozer, P.; Hornick, J.L.; Cabaraux, J.F.
\newblock Forest and silvopastoral cover changes and its drivers in northern
  Morocco.
\newblock {\em Appl. Geogr.} {\bf 2018}, {\em 101},~23--35.

\bibitem[Aouatif~Adghir and Kettani(2018)]{Adghir02112018}
Aouatif~Adghir, H.D.J.; Kettani, K.
\newblock The Tipulidae (Diptera) of northern Morocco with a focus on the Rif
  region, including the description of a new species of Tipula (Lunatipula) and
  an updated checklist for Morocco.
\newblock {\em Ann. SociéTé Entomol. Fr. (N.S.)} {\bf
  2018}, {\em 54},~522--538.
\newblock {\url{https://doi.org/10.1080/00379271.2018.1530949}}.

\bibitem[Boutahar et~al.(2024)Boutahar, Gonzalez, Picone, Crisafulli, Mesa,
  Redouan, Bakali, Kadiri, Lamrani, and Merzouki]{Boutahar01102024}
Boutahar, A.; Gonzalez, P.C.; Picone, R.M.; Crisafulli, A.; Mesa, J.M.;
  Redouan, F.Z.; Bakali, I.E.; Kadiri, M.; Lamrani, Z.; Merzouki, A.
\newblock The relationships between modern pollen rain assemblages and
  vegetation from Sougna Mountain (Rif Mountains-Northern Morocco).
\newblock {\em Palynology} {\bf 2024}, {\em 48},~2358016.
\newblock {\url{https://doi.org/10.1080/01916122.2024.2358016}}.

\bibitem[J.~A.~Jiménez and Guerra(2002)]{Jimenez01092002}
J.~A.~Jiménez, R. M.~Ros, M.J.C.; Guerra, J.
\newblock Contribution to the bryophyte flora of Morocco: terricolous and
  saxicolous bryophytes of the Jbel Bouhalla.
\newblock {\em J. Bryol.} {\bf 2002}, {\em 24},~243--250.
\newblock {\url{https://doi.org/10.1179/037366802125001411}}.

\bibitem[Lahcen~Dahmani and Dindaroglu(2024)]{Dahmani10102024}
Lahcen~Dahmani, Said~Laaribya, H.N.; Dindaroglu, T.
\newblock Landslide hazard mapping in Chefchaouen, Morocco: AHP-GIS
  integration.
\newblock {\em Int. J. Environ. Stud.} {\bf 2024}, {\em
  24},~12483
\newblock {\url{https://doi.org/10.1080/00207233.2024.2412483}}.

\bibitem[Es-Smairi et~al.(2022)Es-Smairi, Moutchou, Touhami, Namous, and
  Mir]{EsSmairi13122022}
Es-Smairi, A.; Moutchou, B.E.; Touhami, A.E.O.; Namous, M.; Mir, R.A.
\newblock Landslide susceptibility mapping using GIS-based bivariate models in
  the Rif chain (northernmost Morocco).
\newblock {\em Geocarto Int.} {\bf 2022}, {\em 37},~15347--15377.
\newblock {\url{https://doi.org/10.1080/10106049.2022.2097322}}.

\bibitem[Manaouch et~al.(2024)Manaouch, Sadiki, Aghad, Pham, Batchi, and
  Karkouri]{Manaouch03032024}
Manaouch, M.; Sadiki, M.; Aghad, M.; Pham, Q.B.; Batchi, M.; Karkouri, J.A.
\newblock Assessment of landslide susceptibility using machine learning
  classifiers in Ziz upper watershed, SE Morocco.
\newblock {\em Phys. Geogr.} {\bf 2024}, {\em 45},~203--230.
\newblock {\url{https://doi.org/10.1080/02723646.2023.2250174}}.

\end{thebibliography}


%%%%%%%%%%%%%%%%%%%%%%%%%%%%%%%%%%%%%%%%%%
\PublishersNote{}
%\isPreprints{}{% This command is only used for ``preprints''.
\end{adjustwidth}
%} % If the paper is ``preprints'', please uncomment this parenthesis.
\end{document}
